\PassOptionsToPackage{hyphens}{url}
\documentclass[aps,prd,twocolumn,superscriptaddress,nofootinbib]{revtex4-2}

% MiKTeX REVTeX 4.2 needs explicit package loading
\usepackage[utf8]{inputenc}
\usepackage[T1]{fontenc}
\usepackage{lmodern}
\usepackage{amsmath}
\usepackage{amssymb}
\usepackage{amsthm}
\usepackage{booktabs}
\usepackage{tabularx}
\usepackage{xcolor}
\usepackage{graphicx}
\graphicspath{{../figures/}}
\input{glyphtounicode}
\pdfgentounicode=1

\newtheorem{definition}{Definition}
\newtheorem{proposition}{Proposition}
\newtheorem{theorem}{Theorem}
\newtheorem{corollary}{Korollar}
\newtheorem{lemma}{Lemma}
\newtheorem{remark}{Bemerkung}
\newtheorem{axiom}{Axiom}
\makeatletter
\def\Dated@name{Stand:\ }
\makeatother
\newcommand{\germanmonthname}{%
  \ifcase\month\or Januar\or Februar\or März\or April\or Mai\or Juni\or Juli\or August\or September\or Oktober\or November\or Dezember\fi%
}
\newcommand{\germantoday}{\number\day.~\germanmonthname~\number\year}

\usepackage{hyperref}
\hypersetup{
    unicode=true,
    pdftitle={Das Sättigungstheorem: Axiomatische Zwangsbedingungen für Quantengravitation aus kosmologischer Relaxation},
    pdfauthor={Lukas Geiger},
    pdfsubject={Axiomatische Bedingungen für eindimensionale kosmologische Sättigungsdynamik in der Quantengravitation},
    pdfkeywords={Quantengravitation, Sättigungsdynamik, axiomatische Kosmologie, Curvature Relaxation Model, Renormierungsgruppe, tanh-Sättigung},
    colorlinks=true,
    linkcolor=black,
    urlcolor=blue,
    citecolor=black
}
\renewcommand{\acknowledgmentsname}{KI-Offenlegung}

\begin{document}
\pagenumbering{gobble}
\onecolumngrid

% ===================================================================
% TITELSEITE
% ===================================================================

\thispagestyle{empty}
\vspace*{0.08\textheight}

\begin{center}
{\LARGE Das Sättigungstheorem: Axiomatische Zwangsbedingungen\par}
\vspace{0.2em}
{\LARGE für Quantengravitation aus kosmologischer Relaxation\par}
\vspace{1.8em}
{\large Lukas Geiger\textsuperscript{1,*}\par}
\vspace{0.6em}
{\itshape \textsuperscript{1}Unabhängiger Forscher, Bernau, Deutschland\textsuperscript{\textdagger}\par}
\vspace{0.2em}
{(Stand: \germantoday)\par}
\end{center}

\vspace{1.8em}
\noindent
Wir formulieren vier minimale Axiome -- endliche geometrische Kapazität, kooperative Verstärkung mit antisymmetrischer Erweiterung, monotone kosmische Kontrolle und Renormierungsgruppen-Komposition -- als makroskopische Kandidatenbedingungen für jede Quantentheorie der Gravitation, deren kosmologischer Limes die allgemeine Relativitätstheorie mit einem beschränkten Relaxationsmechanismus reproduziert. Diese Axiome liefern eine Abel-Koordinatenbeschreibung eindimensionaler Sättigungsdynamik, sobald die reversible eindimensionale Antwortkanal-Hypothese $D^\pm$ hinzugenommen wird. Ein Tangens-hyperbolicus-Profil folgt, sobald der Sättigungsrand zusätzlich durch die projektive Doppelrand-Quotienten-Annahme D' festgelegt wird: der Quotient $R(x)=(1+x)/(1-x)$ muss sich unter Komposition multiplizieren. Ohne diese Annahme lässt glatte Randregularität allein eine Familie glatter hyperbolischer Randkrägen zu, etwa Abel-Koordinaten der Form $\operatorname{arctanh}(x)+\varepsilon x$. Das korrigierte \textit{Sättigungstheorem} ist daher ein projektives Randkragen-Normalformtheorem: $S(g)=S_{\max}\tanh(\alpha\,u(g))$ wird durch Skalenkohärenz plus projektiven Randquotienten ausgewählt, nicht durch bloße Glattheit allein. Große Quantengravitationsprogramme -- Schleifen-Quantengravitation, Asymptotic Safety, String-/holographische Ansätze, Kausalmengentheorie und nichtkommutative Geometrie -- liefern Motivation für A--D, wenn sie auf kosmologische Observablen eingeschränkt werden; D' bleibt dagegen eine separate diagnostische Annahme, die physikalisch erst herzuleiten ist. Unter diesen Hypothesen kann die $\tanh$-Sättigungs-ODE $dX/da = k(1-X^2)$ des Curvature Relaxation Model als projektive Normalform kapazitätsbegrenzter kooperativer Relaxation gelesen werden, wobei offen bleibt, welche UV-Vervollständigung $k$, $\Phi_0$ und den projektiven Quotienten physikalisch festlegt.

\vfill
{\small
\noindent\textsuperscript{*} ORCID: \url{https://orcid.org/0009-0005-7296-1534}\par
\noindent\textsuperscript{\textdagger} Details zum Werkzeugeinsatz stehen in der KI-Offenlegung am Ende des Papers. Alle wissenschaftlichen Inhalte verantwortet ausschließlich der Autor. Paper~V der CRM-Serie. Begleitarbeiten: \cite{Geiger2026,Geiger2026b,Geiger2026c,Geiger2026d}.\par
}

\clearpage
\twocolumngrid
\pagenumbering{arabic}

\iffalse

% ===================================================================
% TITELSEITE
% ===================================================================

\title{Das Sättigungstheorem: Axiomatische Zwangsbedingungen\texorpdfstring{\\}{ }für Quantengravitation aus kosmologischer Relaxation}

\author{Lukas Geiger}
\thanks{ORCID: \url{https://orcid.org/0009-0005-7296-1534}}
\affiliation{Unabhängiger Forscher, Bernau, Deutschland}

\date{\germantoday}

\begin{abstract}
Wir formulieren vier minimale Axiome -- endliche geometrische Kapazität, kooperative Verstärkung mit antisymmetrischer Erweiterung, monotone kosmische Kontrolle und Renormierungsgruppen-Komposition -- als makroskopische Kandidatenbedingungen für jede Quantentheorie der Gravitation, deren kosmologischer Limes die allgemeine Relativitätstheorie mit einem beschränkten Relaxationsmechanismus reproduziert. Diese Axiome liefern eine Abel-Koordinatenbeschreibung eindimensionaler Sättigungsdynamik, sobald die reversible eindimensionale Antwortkanal-Hypothese $D^\pm$ hinzugenommen wird. Ein Tangens-hyperbolicus-Profil folgt, sobald der Sättigungsrand zusätzlich durch die projektive Doppelrand-Quotienten-Annahme D' festgelegt wird: der Quotient $R(x)=(1+x)/(1-x)$ muss sich unter Komposition multiplizieren. Ohne diese Annahme lässt glatte Randregularität allein eine Familie glatter hyperbolischer Randkrägen zu, etwa Abel-Koordinaten der Form $\operatorname{arctanh}(x)+\varepsilon x$. Das korrigierte \textit{Sättigungstheorem} ist daher ein projektives Randkragen-Normalformtheorem: $S(g)=S_{\max}\tanh(\alpha\,u(g))$ wird durch Skalenkohärenz plus projektiven Randquotienten ausgewählt, nicht durch bloße Glattheit allein. Große Quantengravitationsprogramme -- Schleifen-Quantengravitation, Asymptotic Safety, String-/holographische Ansätze, Kausalmengentheorie und nichtkommutative Geometrie -- liefern Motivation für A--D, wenn sie auf kosmologische Observablen eingeschränkt werden; D' bleibt dagegen eine separate diagnostische Annahme, die physikalisch erst herzuleiten ist. Unter diesen zusätzlichen Hypothesen kann die $\tanh$-Sättigungs-ODE $dX/da = k(1-X^2)$ des Curvature Relaxation Model als kanonische projektive Normalform kapazitätsbegrenzter kooperativer Relaxation gelesen werden; offen bleibt, welche UV-Vervollständigung $k$, $\Phi_0$ und den projektiven Quotienten physikalisch festlegt.
\end{abstract}

\keywords{Quantengravitation, Sättigungsdynamik, axiomatische Kosmologie, konditionales Normalformtheorem, Curvature Relaxation Model, tanh-Normalform, Renormierungsgruppe}

\thanks{Details zum Werkzeugeinsatz stehen in der KI-Offenlegung am Ende des Papers. Alle wissenschaftlichen Inhalte verantwortet ausschließlich der Autor.}

\thanks{Paper~V der CRM-Serie. Begleitarbeiten: \cite{Geiger2026,Geiger2026b,Geiger2026c,Geiger2026d}.}

\begin{titlepage}
\maketitle
\end{titlepage}

\fi

% ===================================================================
% 1. EINLEITUNG
% ===================================================================
\section{Einleitung}
\label{sec:intro}

\subsection{Das Problem}

Die Vereinigung der Quantenmechanik und der allgemeinen Relativitätstheorie bleibt das zentrale offene Problem der theoretischen Physik. Nach Jahrzehnten intensiver Forschung haben mehrere bedeutende Forschungsprogramme tiefe strukturelle Einsichten geliefert, ohne bislang auf eine einzige Theorie zu konvergieren:

\begin{itemize}
\item \textbf{Effective field theory (EFT) der Gravitation} berechnet zuverlässige Quantenkorrektionen bei niedrigen Energien, bietet aber keine UV-Vervollständigung \cite{Donoghue1994}.
\item \textbf{Asymptotic Safety} sucht einen UV-Fixpunkt des Renormierungsgruppenflusses (RG) der Gravitation, wodurch die Theorie ohne neue Freiheitsgrade eingeschränkt wird \cite{Weinberg1979,Reuter1998,Percacci2017}.
\item \textbf{Stringtheorie} bietet einen UV-endlichen Rahmen, der Gravitation automatisch enthält, sieht sich aber dem Landscape/Selection-Problem gegenüber \cite{Polchinski1998}.
\item \textbf{Loop Quantum Gravity (LQG)} erreicht Hintergrundunabhängigkeit mit diskreten Geometrieoperatoren, hat aber Schwierigkeiten mit der Dynamik und dem klassischen Limes \cite{Rovelli2004,Thiemann2007}.
\item \textbf{Holographie/AdS-CFT} bietet eine nichtperturbative Definition der Quantengravitation im anti-de~Sitter-Raum, mit unklarer Erweiterung auf realistische Kosmologie \cite{Maldacena1999}.
\item \textbf{Causal Sets} postulieren diskrete Kausalstruktur als fundamental, mit dem Kontinuum als emergent \cite{Sorkin2003,Surya2019}.
\item \textbf{Noncommutative geometry} formuliert Geometrie als Spektraldaten neu und verbindet natürlich Operatoren und Raumzeit \cite{Connes1994}.
\end{itemize}

Orthogonal zu diesen UV-Vervollständigungsprogrammen beschränken die \textit{Swampland-Vermutungen} \cite{Obied2018} die Frage, welche niederenergetischen Theorien konsistent in die Quantengravitation eingebettet werden können, während die \textit{EFT der Dark Energy} \cite{Gubitosi2013} Abweichungen von $\Lambda$CDM auf perturbativer Ebene parametrisiert. Kein Framework bestimmt jedoch die \textit{funktionale Form} des kosmologischen Sättigungsprofils.

Mehrere UV-Vervollständigungsprogramme enthalten Motive, die mit einer beschränkten kosmologischen Antwort verglichen werden können: endliche Skalen, monotone Zweige, Vergröberung oder geordnete Phasen. All diese Motive ``Sättigung'' oder ``Kooperation'' zu nennen, ist zunächst eine Quell-zu-Ziel-Analogie und kein gemeinsamer mikroskopischer Mechanismus. Dieses Papier fragt, welche expliziten Invarianten für eine konditionale Normalformaussage ausreichen.

\subsection{Das CRM als empirischer Anker}

Das Curvature Relaxation Model (CRM) \cite{Geiger2026,Geiger2026b,Geiger2026c,Geiger2026d} bietet einen konkreten, empirisch getesteten kosmologischen Rahmen, in dem Sättigungsdynamik eine zentrale Rolle spielt. Das Modell ersetzt die kosmologische Konstante $\Lambda$ und partikuläre Dunkle Materie durch geometrische Krümmungsrückkehr und erreicht $\Delta\chi^2 = -5,5$ gegenüber $\Lambda$CDM auf gemeinsamen SN+CMB+BAO-Daten mit 6 Parametern und null Early Dark Energy \cite{Geiger2026b}. Das gesamte Framework leitet sich aus einer einzelnen dynamischen Gleichung ab, hier in der additiven Normalform-Koordinate $u$ geschrieben:
\begin{equation}
\frac{dX}{du} = \alpha(1 - X^2)\,,
\label{eq:saturation_ode}
\end{equation}
wobei $X = \Omega_\Phi/\Phi_0$ das normalisierte Krümmungsrückkehrpotential ist. Das phänomenologische CRM verwendet den Skalenfaktor $a$ als Parameter, mit Lösung $X(a) = \tanh(k(a - a_{\mathrm{trans}}))$, die Spätzeit-Beschleunigung (``Dunkle Energie'') in ihrer gesättigten Phase und gravitatives Gerüst (``Dunkle Materie'') in ihrer ungesättigten Phase bietet \cite{Geiger2026b}. \emph{Koordinaten-Caveat (in dieser Revision ergänzt):} Wird ein allgemeiner Parameter $a$ statt der additiven Normalform-Koordinate $u$ verwendet, so gilt $dX/da = \alpha\,u'(a)(1-X^2)$; die CRM-Gleichung mit konstantem Koeffizienten in $a$ erfordert daher die \emph{zusätzliche} affine-Kanal-Annahme $u(a) = \lambda(a - a_{\mathrm{trans}})$. Das Theorem dieses Papers etabliert die projektive Normalform in $u$; es etabliert nicht diese affine Identifikation für irgendeine bestimmte UV-Vervollständigung (siehe Abschnitt~\ref{sec:discussion}).

Paper III \cite{Geiger2026c} schlug diese ODE historisch über eine direkt spurgekoppelte effektive Lagrangedichte vor. Das Wahrheitsaudit vom Juli~2026 bewertete diese Route als unzureichende physikalische Herleitung und behielt sie nur als dokumentiertes Fehlaudit bei. Die aktive Route-3-Rekonstruktion trennt stattdessen einen shiftsymmetrischen Staubkandidaten, einen Vektor--Skalar-MOND-Sektor und den Sättigungsskalar; sie hat weder die Sättigungs-ODE noch D' aus einer gemeinsamen Lagrangedichte hergeleitet. Das nachfolgende Theorem ist daher logisch unabhängig vom historischen Herleitungsclaim.

Das vorliegende Papier ordnet diese Beobachtung als \textit{strukturelle Analogie mit konditionalem Normalformgehalt} ein: die vier physikalisch motivierten Axiome liefern zusammen mit B' und $D^\pm$ ein Abel-Koordinatengesetz der Sättigung, und die zusätzliche projektive Doppelrand-Quotienten-Annahme D' wählt den $\tanh$-Repräsentanten aus.

\subsection{Umfang und Struktur}

Das zentrale Ergebnis ist das \textit{Saturation Theorem} (Theorem~\ref{thm:saturation}): Die Axiome A--D reduzieren zusammen mit der signierten Antisymmetrie B' und der separaten Hypothese $D^\pm$ des eindimensionalen reversiblen Kanals die effektive kosmologische Antwort auf eine Abel-Koordinate, während die projektive Doppelrand-Quotienten-Annahme D' die $\tanh$-Normalform bis auf Umparametrisierung eindeutig auswählt. Das Papier ist wie folgt strukturiert:

\begin{itemize}
\item Section~\ref{sec:axioms}: Formulierung der vier Axiome.
\item Section~\ref{sec:theorem}: Aussage und Beweis des Saturation Theorem.
\item Section~\ref{sec:exclusion}: Nichtprojektive Kompositionsgesetze statt Profilnamen.
\item Section~\ref{sec:qg}: Kompatibilitätskarte mit Quantengravitations-Motiven.
\item Section~\ref{sec:ontology}: Ontologische Implikationen: Raumzeit als geordnete Phase.
\item Section~\ref{sec:observational}: Beobachtungsfolgerungen und offene Fragen.
\item Section~\ref{sec:discussion}: Diskussion und Ausblick.
\end{itemize}


% ===================================================================
% 2. DIE VIER AXIOME
% ===================================================================
\section{Die vier Axiome}
\label{sec:axioms}

Wir formulieren die Axiome in Begriffen einer effektiven Antwortfunktion $S_T$. Um die Rollen sauber zu trennen, bezeichnet $r\geq0$ eine physikalische UV-zu-IR-Kontrollgröße oder ein Skalenverhältnis, $u$ die additive signierte Normalform-Koordinate, die nach Übergang zu einem eindimensionalen Kompositionskanal verwendet wird, und $x=S/S_{\max}$ die normalisierte Antwort. Das Symbol $g$ wird im Folgenden als generische Kanalvariable beibehalten, wenn keine Unterscheidung zwischen $r$ und $u$ nötig ist. Die Funktion $S_T$ kodiert, wie die geometrische Ausgabe der Theorie (Krümmung, Expansionsrate, geometrische Ordnung) auf den Kontrollparameter reagiert.

\subsection{Axiom A: Endliche geometrische Kapazität}

\begin{axiom}[Endliche geometrische Kapazität]
\label{ax:capacity}
Pro Hubble-Patch existiert eine endliche maximale geometrische Ordnung $C_{\max} < \infty$. Äquivalent ist die effektive Antwortfunktion $S_T(r)$ auf dem gewählten positiven physikalischen Antwortkanal global beschränkt:
\begin{equation}
0 \leq S_T(r) < S_{\max} < \infty \quad \forall\, r \geq 0\,.
\end{equation}
\end{axiom}

\textit{Physikalischer Inhalt.} Dieses Axiom fordert endliche Kapazität für den gewählten skalaren Antwortkanal. Es behauptet nicht, dass jede mikroskopische Observable oder UV-Amplitude endlich ist. Endliche Flächenquanta in LQG \cite{Rovelli2004}, die Bekenstein-Schranke \cite{Bekenstein1981}, endliche Codesubräume \cite{Almheiri2015}, Causal-Set-Diskretheit \cite{Sorkin2003} und Fixpunktverhalten \cite{Reuter1998} sind nur Quelldomänenmotive. Jedes benötigt eine explizite Abbildung auf die beschränkte Antwort $S_T$, bevor es als strukturelle Stütze von Axiom~A zählt.

Ohne endliche Kapazität würde der Rahmen unbegrenzte geometrische Konfigurationen zulassen, die mit den singulären Verhaltensweisen assoziiert sind, die Quantengravitationsprogramme typischerweise zu zähmen versuchen.

\textit{Verbindung zum CRM.} Die Sättigungsamplitude $\Phi_0$ im CRM \cite{Geiger2026} ist die physikalische Realisierung von $S_{\max}$. Das Pöschl-Teller-Potential $V(\phi) = V_0/\cosh^2(\phi/\phi_0)$ erzwingt dynamisch $\Omega_\Phi \leq \Phi_0$ \cite{Geiger2026c}.

\subsection{Axiom B: Kooperative Verstärkung}

\begin{axiom}[Kooperative Verstärkung]
\label{ax:cooperative}
Existierende geometrische Ordnung erleichtert weitere Ordnung. Die effektive Antwortfunktion $S_T(g)$ ist glatt ($C^\infty$) für $g \geq 0$ mit einer lokalen Expansion bei $g = 0$:
\begin{equation}
S_T(g) = c_1 g + c_3 g^3 + \mathcal{O}(g^5)\,, \quad c_1 > 0\,.
\label{eq:ir_expansion}
\end{equation}
Die ungerade-Potenz-Struktur spiegelt den katalytischen Charakter der Antwort: die Rate der geometrischen Ordnung ist proportional sowohl zur bestehenden geordneten Fraktion als auch zur verbleibenden Kapazität.
\end{axiom}

\textit{Physikalischer Inhalt.} Formal erhält $c_1>0$ nur das Vorzeichen der lokalen Antwort, während die ungerade Potenzstruktur mit der signierten Erweiterung (Axiom~B') verträglich ist. Das Wort ``kooperativ'' ist eine motivierende Analogie zu positiver Rückkopplung eines Ordnungsparameters; es identifiziert keine mikroskopischen Bestandteile, Wechselwirkungen, Korrelationen oder einen ferromagnetischen Hamiltonoperator.

Innerhalb dieses effektiven Modells würde das Entfernen der positiven lokalen Antwort die beabsichtigte selbstverstärkende Sättigungsinterpretation entfernen. Es würde nicht beweisen, dass klassische Raumzeit unmöglich ist. Spontaner Symmetriebruch wird daher unten als heuristische Quelldomäne verwendet, deren erhaltene und nicht erhaltene Invarianten explizit ausgewiesen werden.

\textit{Verbindung zum CRM.} Die Sättigungskopplung $k$ in der CRM-ODE~\eqref{eq:saturation_ode} liefert den positiven lokalen Koeffizienten $c_1$. Das spieltheoretische Gleichgewicht aus Paper~I \cite{Geiger2026} motiviert die Sprache kooperativer Ordnung; es ist keine mikroskopische Herleitung dieses Mechanismus.

Eine Begleitbedingung erweitert den Bereich von $S_T$ auf negative Argumente, was die gruppentheoretische Klassifizierung von Schritt 2 im Beweis ermöglicht:

\begin{axiom}[Antisymmetrische Erweiterung (B')]
\label{ax:antisymmetry}
Die Antwortfunktion $S_T$ lässt eine glatte antisymmetrische Erweiterung auf alle $g \in \mathbb{R}$ zu:
\begin{equation}
S_T(-g) = -S_T(g) \quad \forall\, g\,.
\label{eq:antisymmetry}
\end{equation}
\end{axiom}

\textit{Physikalischer Inhalt.} Die Antisymmetrie hat drei komplementäre Motivationen:

\textit{(i) Vorzeichenkovarianz der Antwort.} Der Parameter $g$ ist keine ``Entfernung'' (die per Definition nicht-negativ wäre), sondern eine \textit{vorzeichenbehaftete Abweichung} von einem symmetrischen Referenzzustand. Im CRM parametrisiert $g$ die Krümmungsperturbation relativ zum de~Sitter-Hintergrund; positives $g$ entspricht überdichten Perturbationen, die geometrische Ordnung verstärken, negatives $g$ zu unterdichten Regionen, die sie reduzieren. Vorzeichenkovarianz ($S_T(-g) = -S_T(g)$) besagt dann, dass die Antwortfunktion über- und unterdichte Regionen symmetrisch in Magnitude behandelt -- eine natürliche physikalische Anforderung ohne symmetriebrechenden Mechanismus in der Antwort selbst. Wir betonen, dass Axiom B' \textit{nicht} eine Aussage über kosmologische Zeitumkehrsymmetrie ist (die durch die Expansion gebrochen ist); es ist eine Beschränkung der funktionalen Form der Antwort unter $g \to -g$.

\textit{(ii) Präzedenzfälle in der Physik.} Vorzeichenkovarianz von Antwortfunktionen ist in symmetrischen Antwortkanälen vertraut: in der Elektrodynamik die Polarisierung $P(-E) = -P(E)$ für zentrosymmetrische Medien und im Magnetismus $M(-H) = -M(H)$ für die idealisierte paramagnetische/ferromagnetische Antwort. Die ungerade Erweiterung in Axiom~B' ist eine Modellannahme in einer lokal signierten Antwortkoordinate; es wird nicht behauptet, dass generische Beta-Funktionen ungerade sind.

\textit{(iii) Mathematische Rolle.} Zusammen mit der ungerade-Potenz-Struktur von Axiom B erweitert B' das Kompositionsgesetz von der Halbgerade $[0, S_{\max})$ zu einer vollständigen Gruppenoperation auf $(-S_{\max}, S_{\max})$, die für die Abel-Gleichungs-Klassifizierung wesentlich ist. Ohne B' steht nur eine Halbgruppen-Struktur zur Verfügung, und das Abel-Darstellungsargument von Schritt 2 gilt nicht in derselben Form. Wir beachten, dass dies eine explizite Annahme ist, keine abgeleitete Konsequenz; ihre Lockerung führt zu alternativen Kompositionsstrukturen (z.B. $C(x,y) = x + y - xy$ auf $[0,1)$), weshalb sie als separates Axiom erscheint, statt in Axiom B subsumiert zu werden.

\subsection{Axiom C: Monotone kosmische Kontrolle}

\begin{axiom}[Monotone kosmische Kontrolle]
\label{ax:monotone}
Die effektive Antwortfunktion $S_T(g)$ ist monoton wachsend mit Sättigung im Unendlichen:
\begin{equation}
S_T'(g) \geq 0\,, \quad \lim_{g \to \infty} S_T(g) = S_{\max}\,.
\label{eq:monotone}
\end{equation}
Es gibt genau eine relevante Richtung (Kontrollparameter), die den Übergang von niedriger zu hoher geometrischer Ordnung regiert.
\end{axiom}

\textit{Physikalischer Inhalt.} Dies schließt oszillierende, Multi-Fixpunkt- oder chaotische geometrische Dynamik aus. Es garantiert:
\begin{itemize}
\item Eine eindeutige Zeitrichtung (thermodynamischer Pfeil).
\item Kein ``Springen'' zwischen geometrischen Phasen.
\item Stabile Expansion zu einem de~Sitter-Attraktor.
\end{itemize}

Im CRM ist der Kontrollparameter der Skalenfaktor $a$; in RG-Sprache entspricht er der einzelnen relevanten Richtung, die vom UV-Fixpunkt zum IR fließt \cite{Percacci2017}. Axiom C wird verletzt durch zyklische Kosmologien und durch Modelle mit mehreren stabilen Vakua -- beide denen ein eindeutiger thermodynamischer Pfeil fehlt. Wir beachten, dass Axiom C eine Einschränkung auf den \textit{kosmologischen Sektor} mit einem einzelnen dominanten Übergang ist: es behandelt nicht die String-Landschaft (wo Tunneln zwischen metastabilen Vakua zu nicht-monotonen Trajektorien führt) oder stochastische Inflation (wo lokale Fluktuationen die Monotonie brechen). Diese Szenarios betreffen mehrere effektive Kontrollparameter oder stochastische Dynamik und liegen außerhalb des Bereichs des Einfeld-, deterministischen Axiom-Systems. Das Theorem sollte so gelesen werden: A--D liefern zusammen mit B' und der inneren Gruppenform $D^\pm$ eine Abel-/Randkragenklasse für einen einzelnen deterministischen Sättigungskanal; erst D' wählt den projektiven $\tanh$-Repräsentanten aus.

\textit{Verbindung zum CRM.} Die Monotonie von $\Omega_\Phi(a)$ ist eine direkte Konsequenz der ODE~\eqref{eq:saturation_ode}: da $1 - X^2 > 0$ für $|X| < 1$ ist, ist die Lösung streng wachsend.

\begin{remark}[Axiom-Interdependenz]
Lemma~\ref{lem:global_monotonicity} und Corollary~\ref{cor:monotonicity} isolieren die exakte Abhängigkeit: Die Monotonie- und Randsättigungsklauseln von Axiom~C folgen erst, wenn das reversible Antwortgesetz $D^\pm$ mit einem echten globalen Homomorphismus einer einparametrigen Kontrollgruppe kombiniert wird. Das antwortseitige Gruppengesetz allein liefert diese globale Kontrollabbildung nicht. Wir behalten Axiom~C daher im Haupttheorem explizit bei, sowohl zur physikalischen Transparenz als auch für Modelle, in denen nur ein lokaler Chart, eine Halbgruppe, ein stochastischer Kanal oder eine Mehrfeldkontrolle bekannt ist.
\end{remark}

\subsection{Axiom D: Renormalisierungsgruppenkomposition}

\begin{axiom}[Renormalisierungsgruppenkomposition]
\label{ax:composition}
Es existiert eine Skalenkomposition $g \mapsto g_1 \oplus g_2$ (``zwei RG-Schritte in Abfolge'') derart, dass die Antwortfunktion kompositionell kohärent ist:
\begin{equation}
S_T(g_1 \oplus g_2) = \mathcal{C}\!\left(S_T(g_1),\, S_T(g_2)\right)\,,
\label{eq:composition}
\end{equation}
wobei $\mathcal{C}$ eine glatte, kommutative, assoziative Komposition mit neutralem Element $0$ und absorbierendem Rand $S_{\max}$ ist:
\begin{equation}
\mathcal{C}(S, 0) = S\,, \quad \mathcal{C}(S, S_{\max}) = S_{\max}\,.
\label{eq:composition_boundary}
\end{equation}
Überdies erstreckt sich $\mathcal{C}$ glatt (d.h. $C^\infty$) auf das abgeschlossene Quadrat $[0, S_{\max}]^2$.
\end{axiom}

\begin{axiom}[Eindimensionaler reversibler Kanal ($D^\pm$)]
\label{ax:reversible_channel}
Nach Normalisierung durch $S_{\max}$ und Erweiterung durch B' ist die induzierte Operation auf dem signierten Antwortintervall $C:(-1,1)^2\to(-1,1)$ kürzbar und in jedem Argument strikt monoton, und jede Translation $y\mapsto C(x,y)$ ist ein glatter Diffeomorphismus mit Inverser $y\mapsto C(-x,y)$. Alle nachfolgenden Abel-Koordinaten-Schlüsse verwenden diese zusätzliche eindimensionale deterministische Kanal-Hypothese; sie ist keine Konsequenz des bloßen positiven Rand-Halbgruppengesetzes allein.
\end{axiom}

\textit{Notation für Kanalvariablen.} Sei $r\geq0$ die physikalische UV-zu-IR-Kontrollgröße oder das Skalenverhältnis, $u$ die additive signierte Normalform-Koordinate mit $u(g_1\oplus g_2)=u(g_1)+u(g_2)$ und $x=S/S_{\max}$ die normalisierte Antwort. Aussagen über konstante Koeffizienten sind Aussagen in $u$, nicht automatisch in einem extern gewählten physikalischen Parameter. Der positive physikalische Ast wird durch Einschränkung des signierten lokalen Kanals auf die physikalisch realisierte Seite zurückgewonnen.

\begin{remark}[Randkragen-Guardrail]
Die Glattheitsklausel in Axiom~D wird im Folgenden nicht als Eindeutigkeitsbehauptung für $\mathrm{arctanh}$ verwendet. Die Abel-Koordinate
\[
  \phi_\varepsilon(x)=\operatorname{arctanh}(x)+\varepsilon x
\]
induziert eine glatte assoziative Komposition auf dem offenen Intervall und liefert über
\[
  q_\varepsilon(x)=e^{-2\phi_\varepsilon(x)}
  =\frac{1-x}{1+x}e^{-2\varepsilon x}
\]
einen glatten positiven Randkragen. Bloße Glattheit der positiven Sättigungsfläche wählt daher eine hyperbolische Randkragenklasse, aber keine eindeutige projektive Koordinate. Randglattheit wird als Zulässigkeitsbedingung für den Antwortkanal verwendet. Sie wird in diesem Paper nicht aus einem UV-Modell hergeleitet. Die $\tanh$-Normalform wird erst dann ein Theorem, wenn der projektive Randquotient explizit festgelegt wird.
\end{remark}

\begin{axiom}[Projektive Doppelrand-Quotienten-Annahme (D')]
\label{ax:projective_quotient}
Für die normalisierte Antwort $x=S/S_{\max}\in(-1,1)$ ist der projektive Randquotient
\[
  R(x):=\frac{1+x}{1-x}
\]
unter Komposition multiplikativ:
\begin{equation}
  R(C(x,y))=R(x)R(y).
  \label{eq:projective_quotient}
\end{equation}
Äquivalent ist die signierte Koordinate $h(x):=\frac12\log R(x)=\operatorname{arctanh}(x)$, nicht der Quotient selbst, unter der eindimensionalen Antwortkomposition additiv.
\end{axiom}

\textit{Physikalischer Status von D'.} D' ist die projektive Doppelrand-Quotienten-Annahme. Sobald der Quotient $R=(1+x)/(1-x)$ vor der Wahl einer Abel-Koordinate festgelegt ist, folgt das Geschwindigkeitsadditionsgesetz algebraisch. D' wird somit nicht durch die Hypothesen A--D hergeleitet; es ist der zusätzliche projektive Input, dessen physikalische Herleitung offen bleibt. Es setzt fest, dass die physikalisch komponierte Größe der markierte Randkapazitätsquotient ist, dessen Logarithmus einer Log-Odds- oder Rapiditätskoordinate entspricht. Für jedes UV-Kandidatenprogramm ist der diagnostische Defekt
\[
  \Delta_{D'}(x,y):=\log R(C(x,y))-\log R(x)-\log R(y),
\]
der für den exakten projektiven $\tanh$-Repräsentanten verschwinden muss. $\Delta_{D'}=0$ aus einer mikroskopischen Theorie herzuleiten bleibt eine offene physikalische Aufgabe.

\begin{remark}[Markiert-projektives Natürlichkeitsaudit]
Die Punkte $b_-=-1$, $e=0$ und $b_+=1$ bezeichnen die beiden Antwortränder und die neutrale Antwort. Dann ist
\[
  R(x)=
  \frac{(x-b_-)/(b_+-x)}{(e-b_-)/(b_+-e)}
  =\frac{1+x}{1-x}
\]
ein normiertes Vierpunkt-Doppelverhältnis, und $h(x)=\frac12\log R(x)$ ist die signierte Hilbert-Koordinate von $e$ auf dem Intervall \cite{PapadopoulosTroyanov2008,RainioVuorinen2023}. Diese Größe ist invariant, wenn ein projektiver Koordinatenwechsel auf beide Ränder, den Neutralpunkt und die Antwort wirkt. In einer Karte, die nur die beiden Ränder festhält, gilt für einen orientierungserhaltenden projektiven Automorphismus $R(Tx)=kR(x)$; die zusätzliche Bedingung $T(e)=e$ erzwingt $k=1$, während eine Umkehr der Randorientierung $R$ auf $R^{-1}$ abbildet. Damit ist $R$ nur auf einem orientierten, markierten projektiven Antwortintervall mathematisch kanonisch, nicht auf einem unmarkierten glatten Intervall.

Diese Einschränkung ist physikalischer Zusatzinhalt und keine Konsequenz von A--D oder $D^\pm$. Der obige nichtprojektive Kragen besitzt nämlich unter seiner induzierten Komposition $C_\varepsilon$ die exakt multiplikative Abel-Bilanz
\[
  Q_\varepsilon(x):=e^{2\phi_\varepsilon(x)}
  =R(x)e^{2\varepsilon x},
\]
obwohl der feste Doppelrandquotient $R$ im Allgemeinen $\Delta_{D'}\neq0$ hat. Die Existenz \emph{irgendeiner} multiplikativen Bilanz ist daher generisch und wählt D' nicht aus. Eine physikalische Auswahl von $R$ verlangt eine unabhängig definierte Observable mit zwei endlichen unterscheidbaren Randantworten, einer neutralen Identitätsantwort und einem exakten Kompositionsgesetz, unter dem sich die signierten Hilbert-Inkremente addieren. Eine solche quellenseitig hergeleitete Identifikation wird hier nicht etabliert; D' bleibt eine konditionale projektive Gauge-Fixierung und ein empirisches Defekt-Gate, kein hergeleitetes physikalisches Bilanzgesetz.
\end{remark}

\textit{Physikalischer Gehalt von D.} Axiom~D verlangt Selbstkonsistenz der Antwort unter aufeinanderfolgenden Skalentransformationen. Das spiegelt die Halbgruppeneigenschaft von Renormierungsgruppen-Abbildungen \cite{WilsonKogut1974}; ein struktureller RG-Transfer liegt erst vor, wenn die Antwortprojektion den Homomorphismus~\eqref{eq:rg_response_homomorphism} erfüllt.

Das Kompositionsgesetz~\eqref{eq:composition} kodiert die Anforderung, dass die Theorie \textit{skala-kohärent} im folgenden präzisen Sinne sein muss: die Antwortfunktion $S_T$ ist ein Homomorphismus von der additiven Struktur des Kontrollparameters zur Kompositionsalgebra von Antworten; d.h. keine bevorzugte Skala existiert, auf der diese Homomorphismus-Eigenschaft abbricht oder ihren Charakter ändert. Die absorbierende Randbedingung $\mathcal{C}(S, S_{\max}) = S_{\max}$ spiegelt die Tatsache, dass vollständig gesättigte Geometrie nicht weiter geordnet werden kann -- sie ist ein Fixpunkt.

Die $C^\infty$-Erweiterungsanforderung hat eine direkte physikalische Interpretation: sie verlangt, dass das Kompositionsgesetz \textit{auch nahe der Sättigung} wohlverhalten bleibt. In physikalischen Begriffen bedeutet dies, dass der Ansatz zur geometrischen Kapazitätsgrenze glatt ist und keine neuen Divergenzen oder kritischen Exponenten einführt. Drei physikalische Argumente stützen diese Anforderung:

\textit{(a) Keine neue Phasenübergang bei Sättigung.} Eine Komposition, die singuläre Ableitungen am Rand entwickelt, würde einen Phasenübergang \textit{innerhalb} des Sättigungsprozesses selbst signalisieren -- d.h. der Sättigungsmechanismus würde gerade dort zusammenbrechen, wo er am meisten benötigt wird. Die Glattheitsbedingung schließt solch pathologisches Verhalten aus. Dies unterscheidet sich von Phasenübergängen beim \textit{Beginn} der Sättigung (die erlaubt sind und der Übergangsepoche $a_{\mathrm{trans}}$ des CRM entsprechen); die Randregularität betrifft das \textit{vollständig gesättigte} Regime.

\textit{(b) Unterscheidung von kritischen Phänomenen.} Während kritische Exponenten und divergente Suszeptibilitäten Phasenübergänge bei endlicher Temperatur charakterisieren, ist die Sättigungsgrenze $S \to S_{\max}$ des Theorems ein \textit{asymptotisches} Antwortregime ($g \to \infty$), kein kritischer Punkt. Die ferromagnetische Quelldomäne legt einen Vergleich mit $T\to0$ statt $T\to T_c$ nahe, liefert aber kein binäres Kompositionsgesetz für Magnetisierungen. Randglattheit bleibt daher eine explizite Antwortkanal-Annahme und folgt nicht aus ferromagnetischer Thermodynamik.

\textit{(c) Antwortregularität ist ein Modellpostulat.} Generische Korrelationsfunktionen der Quantenfeldtheorie besitzen Kurzdistanzsingularitäten; Vakuumsprache leitet daher keine $C^\infty$-Randkomposition her. Die Deutung vollständiger Sättigung als maximal geordneter Zustand ist heuristisch, und die verlangte Glattheit bleibt eine explizite Hypothese des gewählten effektiven Antwortkanals.

\textit{Verbindung zum CRM.} Für die phänomenologisch gewählte $\tanh$-Antwort definiert die Additionsidentität $\tanh(u_1 + u_2) = (\tanh u_1 + \tanh u_2)/(1 + \tanh u_1 \tanh u_2)$ algebraisch das verlangte Gesetz $\mathcal{C}(x,y) = (x + y)/(1 + xy/S_{\max}^2)$. Diese interne Normalformidentität ist keine thermodynamische oder mikroskopische Herleitung der physikalischen Kompositionsvariable.


% ===================================================================
% 3. THE SATURATION THEOREM
% ===================================================================
\section{Das Saturation Theorem}
\label{sec:theorem}

\subsection{Aussage}

\begin{theorem}[Saturation Theorem, projektive Randkragenform]
\label{thm:saturation}
Sei $T$ eine Theorie, deren effektive Antwortfunktion $S_T(g)$ die Axiome A--D, die Antisymmetrie-Bedingung B' und die separate Hypothese $D^\pm$ des eindimensionalen reversiblen Kanals von Axiom~\ref{ax:reversible_channel} erfüllt. Dann besitzt die normalisierte Antwort $\sigma(g)=S_T(g)/S_{\max}$ eine ungerade Abel-Koordinate $\phi:(-1,1)\to\mathbb{R}$ mit
\[
  \phi(C(x,y))=\phi(x)+\phi(y).
\]
Wenn zusätzlich die projektive Doppelrand-Quotienten-Annahme D' gilt, dann existieren Konstanten $\alpha > 0$, $S_{\max} > 0$ und eine strikt monotone Umparametrisierung $u = u(g)$ mit $u(0) = 0$, so dass
\begin{equation}
S_T(g) = S_{\max}\,\tanh\!\left(\alpha\,u(g)\right)\,.
\label{eq:saturation_form}
\end{equation}
Das Kompositionsgesetz ist dann als relativistische Geschwindigkeitsaddition $\mathcal{C}(x,y) = (x+y)/(1+xy)$ bestimmt, bis auf Reskalierung durch $S_{\max}$. Ohne D' bestimmen diese Voraussetzungen die Abel-/Randkragenstruktur, schließen aber nicht alle glatten nichtprojektiven Randkrägen aus.
\end{theorem}

Tabelle~\ref{tab:proof_status_ledger} trennt die logischen Schichten des korrigierten Resultats. Dadurch wird verhindert, dass die Kompatibilitätsdiskussion in Abschnitt~\ref{sec:qg} und die empirischen Bemerkungen in Abschnitt~\ref{sec:observational} als Beweis des projektiven Quotienten gelesen werden.

\begin{table*}[t]
\caption{Proof-Status-Ledger für das korrigierte Saturation Theorem.}
\label{tab:proof_status_ledger}
\small
\begin{ruledtabular}
\begin{tabular}{@{}llll@{}}
\parbox[t]{0.14\textwidth}{\textbf{Schicht}} & \parbox[t]{0.24\textwidth}{\textbf{Input}} & \parbox[t]{0.22\textwidth}{\textbf{Output}} & \parbox[t]{0.27\textwidth}{\textbf{Status / Nicht-Behauptung}} \\
\parbox[t]{0.14\textwidth}{A--D, B', $D^\pm$} & \parbox[t]{0.24\textwidth}{Beschränkte, kooperative, signierte eindimensionale innere Komposition} & \parbox[t]{0.22\textwidth}{Ungerade Abel-Koordinate und hyperbolische Randkragenklasse} & \parbox[t]{0.27\textwidth}{Theorem-Level, bedingt auf die Ein-Kanal-Hypothesen; noch kein eindeutiges $\tanh$-Profil.} \\[0.6em]
\parbox[t]{0.14\textwidth}{D'} & \parbox[t]{0.24\textwidth}{Projektiver Randquotient $R(C)=R(x)R(y)$} & \parbox[t]{0.22\textwidth}{Kanonischer $\tanh$-Repräsentant und Geschwindigkeitsadditionsgesetz} & \parbox[t]{0.27\textwidth}{Zusätzliche Normalisierung; diagnostisch geführt, hier nicht aus einem UV-Programm hergeleitet.} \\[0.6em]
\parbox[t]{0.14\textwidth}{QG-Programmkarte} & \parbox[t]{0.24\textwidth}{LQG, Asymptotic Safety, Holographie, Causal Sets, CDT und NCG} & \parbox[t]{0.22\textwidth}{Kompatibilitätsmuster für A--D, $D^\pm$ und D'} & \parbox[t]{0.27\textwidth}{Motivation und Filter; keine Verifikation, dass ein Programm alle Inputs erfüllt.} \\[0.6em]
\parbox[t]{0.14\textwidth}{CRM-Datenanker} & \parbox[t]{0.24\textwidth}{Bestehender CRM-Fit und monotone Übergangsevidenz} & \parbox[t]{0.22\textwidth}{Empirische Motivation für Sättigungsdynamik} & \parbox[t]{0.27\textwidth}{Kein unabhängiger Test dieses Theorems, weil das CRM bereits ein $\tanh$-Profil nutzte.} \\[0.6em]
\parbox[t]{0.14\textwidth}{Gestörte Erweiterungen} & \parbox[t]{0.24\textwidth}{Mehrfeld-, stochastische, EFT- oder UV--IR-Matchingdefekte} & \parbox[t]{0.22\textwidth}{Künftiger Residuum-/Gap-Test $r_X/g_X \ll 1$} & \parbox[t]{0.27\textwidth}{Empfohlene Stabilitätserweiterung; nicht Teil des aktuellen Beweises.} \\
\end{tabular}
\end{ruledtabular}
\end{table*}

\subsection{Zentrale Lemmata}

Die Eindeutigkeit von $\mathrm{arctanh}$ beruht auf einem projektiven Quotienten, nicht auf Glattheit allein.

\begin{lemma}[Projektiver Quotient wählt $\mathrm{arctanh}$]\label{lem:projective_arctanh}
Sei $C$ ein stetiges, strikt monotones, kommutatives und assoziatives Gruppengesetz auf $(-1,1)$ mit Neutralelement $0$. Angenommen, $C$ besitzt eine Abel-Koordinate $\phi$,
\[
  \phi(C(x,y))=\phi(x)+\phi(y),
\]
und erfüllt das projektive Quotientengesetz~\eqref{eq:projective_quotient}. Dann existiert $c>0$ derart, dass
\[
  \phi(x)=c\,\operatorname{arctanh}(x),
  \qquad
  C(x,y)=\frac{x+y}{1+xy}.
\]
\end{lemma}

\begin{proof}
Setze $L(x):=\frac12\log((1+x)/(1-x))=\operatorname{arctanh}(x)$. Nach D' gilt
\[
  L(C(x,y))=L(x)+L(y).
\]
Damit ist $F(t):=\phi(\tanh t)$ eine stetige additive Funktion von $t=L(x)$, also $F(t)=ct$ für ein $c>0$. Folglich gilt $\phi(x)=cL(x)=c\,\operatorname{arctanh}(x)$, und Auflösen von $L(C)=L(x)+L(y)$ ergibt $C(x,y)=(x+y)/(1+xy)$.
\end{proof}

\begin{remark}[Warum D' nötig ist]
Die Familie $\phi_\varepsilon(x)=\operatorname{arctanh}(x)+\varepsilon x$ ist für kleines $\varepsilon$ ungerade, glatt und strikt wachsend, divergiert am Rand und induziert über $q_\varepsilon=e^{-2\phi_\varepsilon}$ einen glatten positiven Sättigungskragen. Die alte Behauptung, $C^\infty$-Randregularität allein schließe alle nicht-$\mathrm{arctanh}$-Abel-Koordinaten aus, war daher zu stark. Das korrigierte Theorem verwendet D' als mathematische und physikalische Normalisierung, die diese Rand-Gauge-Freiheit entfernt.
\end{remark}

\subsection{Beweis des Saturation Theorems}

Der Beweis verläuft in drei Schritten.

\textit{Schritt 1: Herleitung der Funktionalgleichung.}
Definiere die normalisierte Antwort $\sigma(g) = S_T(g)/S_{\max} \in [0, 1)$. Die Axiome A und C garantieren, dass $\sigma$ beschränkt und monoton wachsend ist. Axiom D [Gl.~\eqref{eq:composition}] fordert die Existenz einer Komposition:
\begin{equation}
\sigma(g_1 \oplus g_2) = C(\sigma(g_1), \sigma(g_2))\,,
\label{eq:normalized_composition}
\end{equation}
wobei $C: [0,1) \times [0,1) \to [0,1)$ kommutativ, assoziativ, mit Neutralelement $0$ und absorbierendem Rand $1$ ist.

\textit{Schritt 2: Reduktion auf Abels Gleichung.}
Nach Axiom B hat $\sigma$ die lokale Entwicklung~\eqref{eq:ir_expansion} mit $c_1 > 0$. Axiom B' [Gl.~\eqref{eq:antisymmetry}] liefert die ungerade Erweiterung $\sigma(-g)=-\sigma(g)$. Zusammen mit der inneren Gruppenform $D^\pm$ von Axiom~D ergibt dies eine glatte kürzbare Gruppenoperation auf $(-1,1)$, deren absorbierende Ränder bei $\pm 1$ nur als Grenzflächen behandelt werden.

Unter $D^\pm$ definieren die Assoziativitätsbedingung $C(C(x,y),z) = C(x, C(y,z))$ und Kommutativität ein \textit{einparametriges stetiges Gruppengesetz} auf $(-1,1)$ (die Taylorentwicklung am Ursprung definiert ein eindimensionales kommutatives formales Gruppengesetz im Sinne von \cite{Hazewinkel1978}; hier ist die Verbindung, dass jede glatte, assoziative, kommutative binäre Operation auf einer Umgebung von $0 \in \mathbb{R}$ mit Neutralelement $0$ durch Taylorentwicklung ein formales Gruppengesetz über $\mathbb{R}$ ist, und umgekehrt wendet sich Aczéls analytischer Satz auf die konvergente Realisierung dieser formalen Struktur auf dem Intervall $(-1,1)$ an). Wir wenden Aczéls Darstellungssatz \cite{Aczel1966} auf das \textit{offene Intervall} $(-1,1)$ an, wobei $C$ glatt und strikt monoton in jedem Argument ist und das Neutralelement $0$ im Inneren liegt. Die absorbierenden Randbedingungen $C(x, \pm 1) = \pm 1$ sind \textit{nicht} Teil des Definitionsbereichs, auf den Aczéls Satz angewendet wird; sie werden als Grenzverhalten behandelt. Konkret ergibt Aczéls Satz eine stetige, strikt monotone Funktion $\phi: (-1,1) \to \mathbb{R}$ mit $\phi(0) = 0$. Die absorbierenden Randbedingungen implizieren dann $\phi(x) \to +\infty$ für $x \to 1^-$ und $\phi(x) \to -\infty$ für $x \to -1^+$ (da $C(x,y) \to 1$ für $y \to 1$ impliziert, dass $\phi(x) + \phi(y) \to \infty$, was $\phi(y) \to \infty$ erfordert):
\begin{equation}
\phi(C(x,y)) = \phi(x) + \phi(y)\,.
\label{eq:abel}
\end{equation}
Dies ist Abels Funktionalgleichung. \emph{(Beweisschritt in dieser Revision korrigiert; eine frühere Version schrieb $\phi'(0)=c_1$ aus Axiom~B, was ein Kategorienfehler war: $c_1$ ist die Ableitung der unnormalisierten Antwort $S_T$ am Ursprung des Kontrollparameters, nicht der Abel-Koordinate $\phi$ am Antwortursprung.)} Die saubere Konstruktion verläuft über das invariante Vektorfeld: Sei $a(x) = \partial_2 C(x,0)$; unter $D^\pm$ sind Translationen orientierungserhaltende Diffeomorphismen, also $a(x) > 0$ auf $(-1,1)$, und man kann $\phi(x) = \int_0^x dt/a(t)$ definieren. Die Invarianz von $a$ unter dem Gruppengesetz liefert dann $\phi(C(x,y)) = \phi(x) + \phi(y)$, und $\phi$ erbt die $C^\infty$-Regularität von $C$ auf $(-1,1)$ direkt aus dieser Integraldarstellung. Die Funktion $\phi$ ist bis auf eine positive multiplikative Konstante bestimmt, und die Antisymmetrie (Axiom~B') erlaubt, sie ungerade zu wählen. Axiom~B liefert $d\sigma/dg(0)=c_1/S_{\max}>0$; wird eine orientierungserhaltende additive Kontrollkoordinate $u=u(g)$ verwendet, so geht $\kappa:=(d\sigma/du)(0)>0$ nur über $(\phi\circ\sigma)'(0)=\phi'(0)\kappa$ ein und fixiert die positive Gesamtskala $\alpha$. An diesem Punkt liefern die Axiome A--D ein Abel-/Randkragen-Theorem, aber noch nicht $\phi=\mathrm{arctanh}$.

\textit{Schritt 2b: Projektiver Doppelrand-Quotient.}
Wenn zusätzlich die Annahme D' gilt, liefert Lemma~\ref{lem:projective_arctanh} $\phi=c\,\mathrm{arctanh}$. D' ist die projektive Doppelrand-Quotienten-Annahme, keine Konsequenz von A--D: Sobald $R=(1+x)/(1-x)$ vor der Wahl einer Abel-Koordinate festgelegt ist, folgt das Geschwindigkeitsadditionsgesetz algebraisch. Dies ergibt das Kompositionsgesetz:
\begin{equation}
C(x, y) = \frac{x + y}{1 + x\,y}\,,
\label{eq:addition_formula}
\end{equation}
was die relativistische Geschwindigkeitsadditionsformel ist. Die Eindeutigkeit ist daher projektiv: sie ist Eindeutigkeit nach Fixierung des Randquotienten $R=(1+x)/(1-x)$, nicht Eindeutigkeit aus bloßer Glattheit aller Abel-Krägen.

\textit{Schritt 3: Invertierung zu $\tanh$.}
Da $\phi(\sigma(g)) = c\,\mathrm{arctanh}(\sigma(g))$ eine additive Funktion des Kontrollparameters sein muss (d.h., $\phi(\sigma(g)) = \alpha\,u(g)$ für eine Umparametrisierung $u$ und eine Konstante $\alpha > 0$), erhalten wir:
\begin{equation}
\sigma(g) = \tanh(\alpha\,u(g))\,,
\end{equation}
und daher $S_T(g) = S_{\max}\,\tanh(\alpha\,u(g))$, nachdem die positive Konstante $c$ in $\alpha u$ absorbiert wurde. Wenn die Axiome A--D, B', $D^\pm$ und D' exakt gelten, ist das Resultat exakt. Approximative Axiom-Erfüllung (z.B. Zusammenbruch des Kompositionsgesetzes bei sehr kleinem $g$, wo Standardgravitation dominiert) würde Korrektionen einführen; eine formale Stabilitätsanalyse solcher Störungen ist ein offenes Problem. \hfill $\square$

\begin{lemma}[Globale Gruppenmonotonie ohne Profilprämisse]
\label{lem:global_monotonicity}
Sei $C:(-1,1)^2\to(-1,1)$ ein $C^1$-Gruppengesetz mit Neutralelement $0$, und jede Translation $L_x(y):=C(x,y)$ sei ein $C^1$-Diffeomorphismus. Sei
\[
  \sigma:(\mathbb{R},+)\longrightarrow((-1,1),C)
\]
ein globaler $C^1$-Homomorphismus mit $\sigma(0)=0$ und $\kappa:=\sigma'(0)>0$. Definiere $a(x):=\partial_2C(x,0)$. Dann gilt
\begin{equation}
  a(x)>0,
  \qquad
  \sigma'(u)=\kappa\,a(\sigma(u))>0
  \quad\text{für jedes }u\in\mathbb{R}.
  \label{eq:global_monotonicity_generator}
\end{equation}
Außerdem ist $\sigma$ eine streng wachsende Bijektion von $\mathbb{R}$ auf $(-1,1)$, und daher gilt $\lim_{u\to\pm\infty}\sigma(u)=\pm1$.
\end{lemma}

\begin{proof}
Da $L_x$ und seine Inverse $C^1$ sind, verschwindet $a(x)=L_x'(0)$ nirgends. Die Funktion $a$ ist stetig, $(-1,1)$ ist zusammenhängend, und aus $C(0,y)=y$ folgt $a(0)=1$; daher ist $a(x)>0$ auf dem ganzen Intervall. Differenziere die globale Homomorphismusidentität
\[
  \sigma(u+v)=C(\sigma(u),\sigma(v))
\]
bezüglich $v$ bei $v=0$. Dies liefert Gl.~\eqref{eq:global_monotonicity_generator} unmittelbar, sodass strenges Wachstum eine Folgerung und keine Prämisse ist. Weil $\sigma'(u)>0$ ist, ist das Bild eine offene Untergruppe der zusammenhängenden Zielgruppe $((-1,1),C)$; eine zusammenhängende Gruppe besitzt keine echte offene Untergruppe, also ist das Bild ganz $(-1,1)$. Eine streng wachsende Surjektion von $\mathbb{R}$ auf $(-1,1)$ hat die angegebenen Randgrenzwerte. Der Beweis verwendet weder eine Monotonie von $\sigma$ noch eine Abel-Koordinate oder die projektive Annahme D'.
\end{proof}

\begin{corollary}[Exakte Monotonieabhängigkeit für Axiom C]
\label{cor:monotonicity}
Für die normalisierte eindimensionale Antwort liefert Axiom~B $d\sigma/dg(0)=c_1/S_{\max}>0$. Besitzt die Kontrolle eine orientierungserhaltende globale additive Koordinate $u=u(g)$ und ist die Antwort für alle $u,v\in\mathbb{R}$ ein Homomorphismus $\sigma:(\mathbb{R},+)\to((-1,1),C)$, so ist $\kappa=(d\sigma/du)(0)>0$, und Lemma~\ref{lem:global_monotonicity} leitet die Monotonie- und Randsättigungsklauseln von Axiom~C her. Die einzelne relevante Richtung ist bereits in der eindimensionalen globalen Kontrollgruppe kodiert; sie ist kein zusätzliches Ergebnis der Ableitungsrechnung.
\end{corollary}

\begin{remark}
Dies ist ein konditionales Redundanzergebnis für die verstärkte globale Ein-Kanal-Unterklasse. Die antwortseitige Hypothese $D^\pm$, ein lokaler additiver Chart oder eine positive lokale Steigung allein etablieren das globale Homomorphismus-Gate nicht. Das Haupttheorem behält Axiom~C daher explizit, wenn dieses Gate für die physikalische Kontrollvariable nicht bewiesen wurde.
\end{remark}

\begin{remark}[Selberg als hyperbolischer Positivkontrollfall mit Lie-Ursprung]
Für eine kompakte hyperbolische Fläche $M=\Gamma\backslash\mathbb{H}^2\simeq\Gamma\backslash\mathrm{PSL}(2,\mathbb{R})/\mathrm{SO}(2)$ besitzt der Selberg-Spurformelrahmen eine klassische quellenseitige Operatorstruktur: Der Laplace--Casimir-Operator kommutiert mit dem sphärischen bi-$K$-invarianten Faltungskanal \cite{Selberg1956}. Das projektive CRM-Gesetz ist dagegen die kollineare eindimensionale Einstein--Möbius-Addition. Mit $x=\tanh u$ wird die Boost-Untergruppe $B(u)\in\mathrm{SO}^{+}(1,1)$ realisiert, für die $B(u)B(v)=B(u+v)$ und damit $C(x,y)=(x+y)/(1+xy)$ gilt \cite{Ungar2007}. Dies ist ein kontrolliertes gemeinsames hyperbolisches Lie-Motiv vom Rang eins; es ist weder eine Identifikation der Selberg-Fläche mit dem Antwortintervall noch eine Herleitung von D'.

Der stärkere, in einem früheren Selberg-Ideenanhang vorgeschlagene ``E10''-Transfer folgt daraus nicht. Endlicher hyperbolischer Flächeninhalt ist keine endliche skalare Antwortkapazität; Casimir-Zentralität ist keine kooperative Verstärkung; Antisymmetrie einer Lie-Klammer ist keine Ungeradheit einer Antwort; und die Komposition von Geodätenfluss oder Faltung liefert nicht den von Axiom~D verlangten Antwort-Homomorphismus. Die Selberg-Konstruktion definiert weder eine monotone kosmologische Kontrolle noch markierte Sättigungsränder oder eine projektive Antwortobservable. Selberg ist daher nur für die Verfügbarkeit eines invarianten Operator-Kanals mit Lie-Ursprung ein Positivkontrollfall. Der Transfer liefert keine Kosmologieaussage, keine UV-Vervollständigung, keine vollständige Verifikation von A--D/B' und keine Verstärkung des Sättigungstheorems.
\end{remark}

\begin{proposition}[Nicht-Redundanz der Axiome]
\label{prop:independence}
Die folgenden Beispiele sind \emph{Trennungsskizzen}, kein vollständiges Unabhängigkeitstheorem; vollständige paarweise Unabhängigkeit des Axiomensatzes wird nicht behauptet. \emph{(Zwei Beispiele in dieser Revision nach externem Review korrigiert.)} Da Axiom~D im beschränkten Antwortraum formuliert ist, ist Punkt (1) als algebraisches Beispiel für die innere Komposition nach Entfernen der Kapazitätsschranke zu lesen, nicht als vollständiges Modell des endlich-kapazitiven Axiomensystems.

\textit{(1) Verletzt A, erfüllt B/B' und den algebraischen inneren Kompositionsteil von D:} $\sigma(g) = g$ (unbegrenzte lineare Antwort; $C(x,y) = x+y$, assoziativ und kommutativ).

\textit{(2) Verletzt die lokale Ungeradheit von B/B', erfüllt A,C,D:} $\sigma(g) = 1 - e^{-g}$ für $g \geq 0$, mit $C(x,y) = x + y - xy$ (probabilistisches Oder). \emph{Korrektur:} Eine frühere Version führte dies als ``verletzt B$'$, erfüllt B'' auf; tatsächlich enthält die Taylorentwicklung $1-e^{-g} = g - g^2/2 + g^3/6 - \cdots$ einen quadratischen Term, sodass das Beispiel die von Axiom~B selbst geforderte ungerade lokale Entwicklung verletzt, nicht nur die signierte Erweiterung B$'$. Es trennt $\{$B,B$'\}$ von $\{$A,C,D$\}$, trennt aber nicht B von B$'$.

\textit{(3) Verletzt C (und A auf dem positiven Ast), erfüllt B/B' lokal, aber nicht $D^\pm$:} $\sigma(g) = \tanh(g) \cdot \cos(\epsilon g)$ für kleine $\epsilon > 0$. \emph{Korrektur:} Eine frühere Version behauptete, dieses Beispiel ``erfüllt A''; tatsächlich wechselt für $g \gg 1$ der Faktor $\cos(\epsilon g)$ das Vorzeichen, sodass das Beispiel auf dem positiven Ast auch die Nichtnegativität von Axiom~A verletzt und keinen Sättigungsgrenzwert im Sinne von Axiom~C besitzt. Es trennt daher $\{$C$\}$ (und monotone Sättigung allgemein) von lokalem $\{$B,B$'\}$, ist aber \emph{kein} Zeuge dafür, dass C unabhängig von A ist. Ein saubererer Kandidatenzeuge für C-Unabhängigkeit -- ein glattes, ungerades, beschränktes Profil der Form $\sigma(g) = \tanh(g) + \varepsilon\,\psi(g)$ mit $\psi$ einer kompakt getragenen ungeraden Bump-Funktion, gewählt so, dass $0 \leq \sigma < 1$ auf $g \geq 0$, aber $\sigma' < 0$ auf einem Intervall -- sollte vor der Behauptung in einem Trennungs-Ledger verifiziert werden. Nach Korollar~\ref{cor:monotonicity} ist C in jedem Fall \textit{nicht} unabhängig von $\{A, B, B', D^\pm\}$ gemeinsam -- es folgt unter dem inneren Gruppengesetz als abgeleitete Konsequenz. C bleibt für physikalische Klarheit und zur expliziten Markierung des Ein-Kanal-Pfeils erhalten.

\textit{(4) Verletzt D, erfüllt A--C,B':} $\sigma(g) = (2/\pi)\arctan(g)$; beschränkt, monoton, antisymmetrisch, aber seine induzierte Komposition (über $\phi = \tan(\pi x/2)$) entwickelt divergente zweite Ableitungen ($\sim \epsilon^{-1}$) am Rand von $[-1,1]^2$, was die $C^\infty$-Erweiterungs-Anforderung von Axiom D verletzt. Ähnlich hat $\sigma(g) = g/\sqrt{1+g^2}$ (mit $\phi(x) = x/\sqrt{1-x^2}$) divergente zweite Ableitungen am Rand ($\sim \epsilon^{-1/2}$).
\end{proposition}

\subsection{Interpretation}

Der Satz sagt aus, dass $\tanh$ keine bloße \textit{Modellwahl} ist, sondern der \textit{kanonische projektive Repräsentant} eines eindimensionalen Abel-Sättigungskanals. Der mathematische Schlüsseleinblick ist zweistufig: Axiom D erzwingt, dass die Antwortfunktion die Inverse einer Abel-Koordinate ist, und D' fixiert den sonst freien Randkragen dadurch, dass der projektive Quotient $R=(1+x)/(1-x)$ multiplikativ ist.

Physikalisch sagt der Satz aus:
\begin{quote}
\textit{Jede Quantengravitationstheorie, die skalenkohärent, kooperativ, beschränkt, monoton ist und deren Sättigungsrand durch den projektiven Quotienten normalisiert wird, muss kosmologische Sättigungsdynamik vom $\tanh$-Typ erzeugen. Ohne diese Quotienten-Normalisierung liefert das Theorem eine Abel-/Randkragenklasse statt eines eindeutigen Profils.}
\end{quote}

Die Umparametrisierung $u(g)$ absorbiert die mikroskopischen Details der Theorie: verschiedene UV-Vervollständigungen entsprechen verschiedenen Funktionen $u(g)$. Die gemeinsame Form $\tanh(\alpha\,u)$ wird erst erzwungen, wenn das projektive Randverhältnis als physikalische Kompositionsvariable identifiziert wurde.


% ===================================================================

% ===================================================================
% 4. AUSSCHLUSS ALTERNATIVER PROFILE
% ===================================================================
\section{Nichtprojektive Kompositionsgesetze}
\label{sec:exclusion}

\begin{corollary}[Nichtprojektive Kompositionsgesetze statt Profilnamen]
\label{cor:exclusion}
Die folgenden Ausschlüsse sind Ausschlüsse der induzierten Antwortkomposition und des Randquotienten unter der gewählten additiven Koordinate. Sie sind keine koordinatenfreien Ausschlüsse von Graphenformen unter beliebiger monotoner Reparametrisierung. Keines der folgenden induzierten Gesetze kann die strukturellen Voraussetzungen des Theorems plus die projektive Quotienten-Annahme D' gleichzeitig erfüllen, außer wenn es durch zulässige Reparametrisierung $u(g)$ auf die $\tanh$-Form reduzierbar ist:
\begin{enumerate}
\item[(i)] \textbf{Logistisches Profil:} $S_{\max}/(1 + e^{-\beta u})$.
\item[(ii)] \textbf{Rationales Profil:} $S_{\max}\,u/(1 + u)$.
\item[(iii)] \textbf{Arkustangens-Profil:} $S_{\max}\,(2/\pi)\arctan(\beta u)$.
\item[(iv)] \textbf{Harte Sättigung:} $S_{\max}\,\min(u, 1)$.
\end{enumerate}
\end{corollary}

\textit{Beweis.} Jedes Alternative verletzt ein spezifisches Axiom:

\textit{(i) Logistisch.} Die logistische Funktion $L(u) = 1/(1 + e^{-\beta u})$ hat $L(0) = 1/2 \neq 0$, was die Bedingung des neutralen Elements verletzt. Die verschobene Logistik $\tilde{L}(u) = (L(u) - 1/2) \cdot 2$ ergibt $\tilde{L}(u) = \tanh(\beta u/2)$ -- genau die $\tanh$-Form. Daher ist jedes ``logistische'' Modell, das Axiom~D erfüllt, in Wirklichkeit $\tanh$ in versteckter Form.

\textit{(ii) Rational.} Die Funktion $R(u) = u/(1+u)$ hat Abel-Transformation $\phi(x) = x/(1-x)$, die $C(x,y) = \phi^{-1}(\phi(x) + \phi(y)) = (x + y - 2xy)/(1 - xy)$ ergibt. Für kleine $x, y$ erhält man $C(x,y) \approx x + y - 2xy$. Der quadratische Mischterm $-2xy$ ist unverträglich mit der Antisymmetriebedingung (Axiom~B'): Antisymmetrisches $\sigma$ impliziert, dass die induzierte Komposition $C(-x,-y) = -C(x,y)$ erfüllen muss, aber $(x+y-2xy)/(1-xy)$ tut das nicht. Äquivalent: $R(u)$ selbst ist nicht antisymmetrisch, was Axiom~B' verletzt.

\textit{(iii) Arkustangens.} Die Funktion $A(u) = (2/\pi)\arctan(\beta u)$ hat Abel-Transformation $\phi(x) = \tan(\pi x/2)$, die $C(x,y) = (2/\pi)\arctan(\tan(\pi x/2) + \tan(\pi y/2))$ ergibt. Während diese Komposition auf dem offenen Quadrat $(-1,1)^2$ wohldefiniert und glatt ist, kann sie nicht glatt auf das geschlossene Quadrat $[-1,1]^2$ erweitert werden: die zweiten partiellen Ableitungen $\partial^2 C/\partial x^2$ und $\partial^2 C/\partial x\,\partial y$ divergieren beide, wenn $(x,y) \to (1,1)$ (numerisch $\sim \epsilon^{-1}$ entlang $x = y = 1 - \epsilon$), was die $C^\infty$-Randregularitätsanforderung von Axiom~D verletzt. Diese algebraische ($\epsilon^{-1}$) Divergenz ist charakteristisch für Abel-Funktionen mit Potenzgesetz-Randverhalten. Logarithmische Krägen sind die zulässige Randklasse an dieser Fläche; D' wählt daraus den projektiven logarithmischen Kragen $\mathrm{arctanh}$ aus.

\textit{(iv) Harte Sättigung.} $\min(u, 1)$ ist bei $u = 1$ nicht $C^\infty$, was die Glattheitsanforderung von Axiom~B verletzt. \hfill $\square$

\begin{remark}
Der logistische Fall (i) ist instruktiv: er scheint eine echte Alternative zu sein, aber Kompositionsgesetz plus projektiver Quotient zwingen ihn, sich auf $\tanh$ zu reduzieren. Dies veranschaulicht die Kraft von Skalenkohärenz zusammen mit D': oberflächliche Freiheit bei der Profilwahl wird auf eine Randkragenwahl reduziert, und der projektive Quotient fixiert den CRM-Rapiditätsrepräsentanten.
\end{remark}


% ===================================================================
% 5. VERIFIKATION IN QUANTENGRAVITATIONS-PROGRAMMEN
% ===================================================================
\section{Kompatibilitätskarte mit Quantengravitations-Motiven}
\label{sec:qg}

Wir kartieren nun, wie A--D natürlicherweise mit Motiven in den großen Quantengravitationsprogrammen kompatibel sind, wenn man sich auf kosmologische Observablen beschränkt. Die stärkere innere Form $D^\pm$ und die projektive Doppelrand-Quotienten-Annahme D' werden separat als offene Diagnostika geführt.

\subsection{Loop Quantum Gravity / Loop Quantum Cosmology}

In Loop Quantum Cosmology (LQC) \cite{Ashtekar2006,Bojowald2008} lautet die modifizierte Friedmann-Gleichung:
\begin{equation}
H^2 = \frac{8\pi G}{3}\,\rho\!\left(1 - \frac{\rho}{\rho_c}\right)\,,
\label{eq:lqc_friedmann}
\end{equation}
wobei $\rho_c \sim \rho_{\mathrm{Pl}}$ eine kritische Dichte ist, die durch den Barbero-Immirzi-Parameter $\gamma_{\mathrm{BI}}$ gesetzt wird. Die Gleichung besitzt eine beschränkte quadratische Antwort, also ein Sättigungsmotiv; sie ist weder die CRM-Antwort-ODE noch ein gemeinsames mikroskopisches Kapazitätsgesetz.

\textit{Axiom A:} Die kritische Dichte $\rho_c$ ist ein Motiv für eine endliche Schranke; ihre Identifikation mit der Kapazität der gewählten CRM-Antwort benötigt jedoch eine Abbildung. \textit{Axiom B:} Die Wiederherstellung der Standard-Friedmann-Gleichung bei $\rho\ll\rho_c$ etabliert weder kooperative Verstärkung noch B'. \textit{Axiom C:} Der expandierende Zweig liefert ein monotones Segment. \textit{Axiom D:} LQC besitzt kein manifestes RG-Kompositionsgesetz im Sinne von Axiom~D; ein rigoroser Antwort-Homomorphismus von der LQG-Dynamik zu Gleichung~\eqref{eq:composition} bleibt offen.

Mit $X = 1 - 2\rho/\rho_c$ kann die LQC-Friedmann-Gleichung~\eqref{eq:lqc_friedmann} als $H^2 \propto (1 - X^2)$ geschrieben werden. Dadurch bleiben ein Polynomfaktor und eine endliche Schranke erhalten, doch die LQC-Gleichung regiert $H^2$, während die CRM-ODE $dX/da$ regiert; ein Homomorphismus zwischen den Dynamiken wird nicht geliefert. Eine vorgeschlagene Parameterkorrespondenz zwischen LQG-Flächenlücke $\Delta$, $\gamma_{\mathrm{BI}}$, $k$ und $\Phi_0$ bleibt daher heuristisch.

\subsection{Asymptotic Safety}

Im Asymptotic Safety-Programm \cite{Weinberg1979,Reuter1998,Percacci2017} hat die gravitationale effektive Wirkung einen UV-Fixpunkt $g^* > 0$, wo die dimensionslose Newton-Konstante und kosmologische Konstante skalenunabhängig werden. Der RG-Fluss von UV zu IR definiert die effektive Gravitationskopplung $G_k$ und kosmologische ``Konstante'' $\Lambda_k$ als Funktionen der RG-Skala $k$.

\textit{Axiom A:} Der UV-Fixpunkt ist ein Motiv für eine endliche Skala, nicht automatisch eine globale Schranke der gewählten kosmologischen Antwort. \textit{Axiom B:} Ein positiver kritischer Exponent definiert eine relevante Richtung; er etabliert keine mikroskopische Kooperation. \textit{Axiom C:} Relevanz allein beweist keine globale Monotonie einer projizierten Observable. \textit{Axiom D:} Die Wilson-RG-Halbgruppe $R_{k_1}\circ R_{k_2}=R_{k_1k_2}$ operiert auf dem Theorieraum \cite{WilsonKogut1974}, nicht direkt auf $S_T$. Ein struktureller Transfer benötigt eine explizite Abbildung $\Pi$ mit
\begin{equation}
\Pi(R_{k_1}\circ R_{k_2})=\mathcal{C}\!\left(\Pi(R_{k_1}),\Pi(R_{k_2})\right).
\label{eq:rg_response_homomorphism}
\end{equation}
Die Beschränkung auf eine dominante relevante Richtung genügt nicht: Ein invarianter eindimensionaler Kanal und der Homomorphismus~\eqref{eq:rg_response_homomorphism} müssen zusätzlich gezeigt werden. Beides wird hier für kein externes UV-Programm hergeleitet.

Asymptotic Safety liefert damit eines der engsten Quelldomänenmotive für Axiom~D, während der Homomorphismus zur Zielantwort offen bleibt.\footnote{Der selbstähnliche, skalenrekursive Charakter von Axiom~D ist nicht
auf den kosmologischen RG-Fluss beschränkt; dasselbe Motiv tritt in der
Verteilung der Primzahlen auf. Im Primonen- bzw.\ Riemann-Gas von
Julia~\cite{Julia1990} besitzt ein System nicht-wechselwirkender bosonischer
Oszillatoren mit Energien $\log p$ über den Primzahlen $p$ die Riemannsche
Zeta-Funktion als Zustandssumme. Hartnoll und Yang~\cite{HartnollYang2025}
haben gezeigt, dass die Belinsky-Khalatnikov-Lifshitz-Dynamik der Gravitation
nahe einer raumartigen Singularität nach semiklassischer Quantisierung eine
modulinvariante konforme Quantenmechanik ist, deren Zustände genau ein solches
primzahlgelabeltes Oszillatorgas realisieren -- mit Wellenfunktionen
proportional zu automorphen $L$-Funktionen auf der kritischen Achse. Dies legt
nahe, dass Schwarzloch-Singularitäten ein weiteres physikalisches System sein
könnten, in dem eine kapazitätsbegrenzte, selbstähnliche geometrische Ordnung
der durch die Axiome kodierten Art natürlicherweise entsteht. Wir verzeichnen
den Zusammenhang als motivierende Parallele, nicht als unabhängige Verifikation
der Axiome~A--D.} Das Asymptotic-Safety-Szenario für Kosmologie \cite{Bonanno2002} produziert beschränktes Verhalten bei hoher Krümmung. Das ist ein Sättigungsmotiv, kein Beleg für die Paper-V-$\tanh$-Antwort oder D'.

Neuere Berechnungen in zunehmend anspruchsvolleren Truncations liefern numerische Evidenz für den Reuter-Fixpunkt. Eichhorn und Held~\cite{EichhornHeld2018} berichten eine Einschränkung der Top-Quark-Yukawa-Kopplung, und Eichhorn~\cite{Eichhorn2020} diskutiert Fixpunktbeschränkungen für Standardmodell-Materie. Diese Resultate stützen die RG-Quelldomäne, verifizieren aber weder skalare Antwortkapazität, kooperative Rückkopplung und globale Monotonie noch den von Axiomen~A--D benötigten Homomorphismus~\eqref{eq:rg_response_homomorphism}.

\subsection{Stringtheorie und Holographie}

In der Stringtheorie motivieren UV-endliche Amplituden eine Analogie endlichen Verhaltens, liefern aber keine endliche Kapazität des gewählten skalaren Antwortkanals. Eine niederenergetische EFT innerhalb einer Kompaktifizierung besitzt einen RG-Fluss; daraus folgen weder Axiome~B--D noch der Antwort-Homomorphismus~\eqref{eq:rg_response_homomorphism}.

Der holographische Zugang (AdS/CFT \cite{Maldacena1999}) verbindet Beschreibungen über Skalen hinweg, während die Ryu--Takayanagi-Formel \cite{Ryu2006} Verschränkungsentropie und geometrische Fläche verknüpft. Das sind Quelldomänenmotive für Skalenbezug und Monotonie, aber kein explizites binäres Kompositionsgesetz auf der Paper-V-Antwort.

Für kosmologische Anwendungen liefern die de-Sitter-Entropieschranke $S_{\mathrm{dS}} = \pi/(\Lambda G) < \infty$ und der holographische RG-Fluss \cite{deBoer2000} Motive für endliche Schranken und Skalenflüsse. Die Abbildung auf $S_T$ und $\mathcal C$ bleibt offen.

\subsection{Causal Set Theory}

In der Causal Set Theory \cite{Sorkin2003,Surya2019} tritt die kosmologische Konstante als dynamische Größe $\Lambda \sim 1/\sqrt{N}$ auf, wobei $N$ die Anzahl der Causal-Set-Elemente ist. Ihre monotone Abnahme ist ein mögliches Kontrollkoordinatenmotiv; sie ist nicht bereits die von Axiom~C verlangte wachsende beschränkte Antwort.

\textit{Axiom A:} Endliche Sprinkling-Dichte ist ein Diskretheitsmotiv, noch nicht die Kapazität einer skalaren Antwort. \textit{Axiom B:} Poisson-Sprinkling ist unkorreliert; die benötigte Invariante positiver Rückkopplung wurde aus sequentiellem Wachstum \cite{Surya2019} nicht hergeleitet. \textit{Axiom C:} Die $1/\sqrt{N}$-Skalierung ist monoton, doch die Umkehr ihrer Orientierung auf $S_T$ erfordert eine gewählte Umparametrisierung. \textit{Axiom D:} Ein Homomorphismus von sequentiellem Wachstum zur Antwortkomposition ist nicht etabliert.

\subsection{Causal Dynamical Triangulations}

Causal Dynamical Triangulations (CDT) \cite{Ambjorn2005,Loll2020} konstruieren Quantengravitation mittels eines regularisierten Pfadintegrals über kausale Triangulationen. Monte-Carlo-Simulationen enthüllen ein Phasendiagramm mit einer physikalisch relevanten ``de Sitter-Phase'' (Phase~$C_{dS}$), in der die emergente Geometrie Hausdorff-Dimension $d_H \approx 4$ hat und ein Volumenprofil, das gut durch euklidischen de Sitter-Raum beschrieben wird.

\textit{Axiom A:} Eine endliche Triangulation ist ein regulatorisches endliches Zählmotiv; eine kontinuierliche skalare Kapazität wurde nicht hergeleitet. \textit{Axiom B:} Wachstum des de-Sitter-Volumenprofils etabliert weder kooperative Rückkopplung noch B'. \textit{Axiom C:} Der expandierende Zweig liefert ein monotones Segment. \textit{Axiom D:} Blocktransformationen sind ein Vergröberungsmotiv, aber ein assoziativer Antwort-Homomorphismus ist nicht etabliert.

\subsection{Nichtkommutative Geometrie}

Das Programm von Connes \cite{Connes1994} rekonstruiert Geometrie aus spektralen Daten (der Dirac-Operator auf einer nichtkommutativen Algebra). Das Spektral-Wirkungs-Prinzip liefert das Standard-Modell-Lagrangian gekoppelt an Gravitation aus dem Spektrum eines einzelnen Operators.

\textit{Axiom A:} Die Spektralabschneidung $\Lambda$ ist nicht bereits eine endliche Kapazität des gewählten Antwortkanals. \textit{Axiom B:} Aus der Wärmekernentwicklung wurde kein signierter ungerader Antwortkanal hergeleitet. \textit{Axiom C:} Eine monotone skalare Projektion ist nicht spezifiziert. \textit{Axiom D:} Produkte spektraler Tripel definieren nicht von selbst das binäre Antwortgesetz $\mathcal C$. Dies bleiben Quelldomänenmotive statt struktureller Transfers.

\subsection{Zusammenfassung: Warum bedingte Abel-/Randkragen-Universalität?}

Tabelle~\ref{tab:axiom_check} fasst die Kompatibilitätskarte zusammen. Für kein einzelnes externes QG-Programm wird hier rigoros bewiesen, dass es A--D, $D^\pm$ und D' gleichzeitig erfüllt. Die Einträge verzeichnen Quelldomänenmotive und offene Abbildungspflichten, keine Transferbeweise. Das Theorem etabliert Universalität nur \emph{innerhalb} der explizit definierten Klasse eindimensionaler Antworten, die seine Voraussetzungen erfüllen: A--D, B' und $D^\pm$ liefern eine Abel-/Randkragenklasse, D' wählt deren projektiven $\tanh$-Repräsentanten. Diese konditionale Normalformäquivalenz ist keine Wilsonsche Universalitätsklasse und impliziert weder gemeinsame mikroskopische Dynamik noch gemeinsame kritische Exponenten \cite{WilsonKogut1974}.

\begin{table*}[t]
\caption{Kompatibilitätskarte mit Quantengravitations-Motiven. $\checkmark$ = kompatibles Motiv mit expliziter Konstruktion, ($\checkmark$) = kompatibles Motiv mit milden Annahmen, ? = offene Herleitung, no = bekannte Obstruktion. Die D'-Spalte notiert, ob der projektive Doppelrand-Quotient hergeleitet oder explizit angenommen wird.}
\label{tab:axiom_check}
\begin{ruledtabular}
\begin{tabular}{lccccc}
Programm & A & B/B' & C & D & D' \\
\hline
LQG/LQC & ($\checkmark$) & ? & ($\checkmark$) & ? & ? \\
Asymptotic Safety & ($\checkmark$) & ? & ? & ($\checkmark$) & ? \\
Strings/Holographie & ($\checkmark$) & ? & ? & ? & ? \\
Causal Sets & ($\checkmark$) & ? & ($\checkmark$) & ? & ? \\
CDT & ($\checkmark$) & ? & ($\checkmark$) & ? & ? \\
NCG & ? & ? & ? & ? & ? \\
CRM (effektiv) & $\checkmark$ & $\checkmark$ & $\checkmark$ & $\checkmark$ & angenommen \\
\end{tabular}
\end{ruledtabular}
\end{table*}

% ===================================================================
% 6. ONTOLOGIE: ORDNUNG, NICHT METRIK
% ===================================================================
\section{Raumzeit als geordnete Phase}
\label{sec:ontology}

Die Kompatibilitätsmotive in Tabelle~\ref{tab:axiom_check} werfen eine konzeptionelle Frage auf: Falls für einen externen Antwortkanal die genannten Gates später tatsächlich gezeigt werden, wie ist diese konditionale Normalformäquivalenz zu interpretieren, ohne eine gemeinsame mikroskopische Ontologie zu folgern?

\subsection{Die konzeptionelle Verschiebung}

Das Saturation Theorem, kombiniert mit dem empirischen Erfolg des CRM, deutet auf eine mögliche interpretive Lesart der Quantengravitation hin:

\begin{quote}
\textit{In dieser Lesart ist Quantengravitation nicht nur die Quantisierung der Metrik, sondern auch die Theorie, wie geometrische Ordnung aus einem vorgemetrischen Quantensystem hervorgeht.}
\end{quote}

Die $\tanh$-Sättigung ist daher am besten als bedingte Signatur eines kapazitätsbegrenzten kooperativen Ordnungsprozesses zu lesen, sobald das eindimensionale Kompositionsgesetz und D' festgelegt sind. Das ist kein unabhängiger Beleg für eine bestimmte mikroskopische Ontologie; es liefert eine gemeinsame Normalformsprache für Spin-Netzwerk-, String-/holographische, Causal-Set- oder Code-Subspace-Realisierungen, sofern sie die genannten Gates erfüllen.

\subsection{Die ferromagnetische Analogie}

Die Analogie zur spontanen Magnetisierung ist auf der Ebene eines Mean-Field-Ginzburg-Landau-Modells nützlich. Die CRM-Sättigungs-ODE $dX/da = k(1 - X^2)$ kann als Bewegungsgleichung für den Ordnungsparameter eines toy-modellhaften Phasenübergangs zweiter Ordnung mit Doppelmulden-Freien-Energie $F(X) = -k(X - X^3/3)$ geschrieben werden. Paper~III \cite{Geiger2026c} entwickelt diese Analogie im Detail und identifiziert:

\begin{center}
\begin{tabular}{ll}
\hline
\shortstack[l]{Ferromagnetische\\Quelldomäne} & \shortstack[l]{CRM-\\Zieldomäne} \\
\hline
Magnetisierung $M$ & Krümmungsrückkehr $\Omega_\Phi$ \\
\shortstack[l]{Inverse Temperatur\\$J/k_BT$} & \shortstack[l]{Kontrollkoordinate\\$a-a_{\mathrm{trans}}$\\(heuristische\\Umparametrisierung)} \\
Curie-Punkt ($J/k_BT_c$) & Übergang $a_{\mathrm{trans}}$ \\
Spin-Wechselwirkung $J$ & Krümmungskopplung $k$ \\
Sättigung $M_s$ & Sättigung $\Phi_0$ \\
$\tanh(J/k_BT)$ & $\tanh(k(a - a_{\mathrm{trans}}))$ \\
\hline
\end{tabular}
\end{center}

\begin{table*}[t]
\caption{Invariantenaudit der in Paper~V verwendeten Quell-zu-Ziel-Transfers. ``Strukturell'' bleibt einer expliziten Normalformabbildung oder einem Homomorphismus vorbehalten; ein gemeinsames Wort oder eine gemeinsame Graphform genügt nicht.}
\label{tab:metaphor_invariants}
\begin{ruledtabular}
\begin{tabular}{llll}
\parbox[t]{0.14\textwidth}{Transfer} & \parbox[t]{0.25\textwidth}{Erhaltene Invariante} & \parbox[t]{0.32\textwidth}{Nicht erhalten oder nicht hergeleitet} & \parbox[t]{0.18\textwidth}{Zulässiger Status} \\
\hline
\parbox[t]{0.14\textwidth}{Sättigung} & \parbox[t]{0.25\textwidth}{Normalisiertes beschränktes Intervall, geordnete Endpunkte, monotone Randannäherung} & \parbox[t]{0.32\textwidth}{Gemeinsame mikroskopische Gleichung, freie Energie oder UV-Mechanismus} & \parbox[t]{0.18\textwidth}{Strukturell nur auf Antwort-/Normalformebene} \\[0.5em]
\parbox[t]{0.14\textwidth}{Ferromagnetisch} & \parbox[t]{0.25\textwidth}{Signierter skalarer Ordnungsparameter und Mean-Field-$\tanh$-förmige Antwort} & \parbox[t]{0.32\textwidth}{Magnetisierungs-Kompositionsgesetz, Gitterkorrelationen, kritische Exponenten, Bestandteile} & \parbox[t]{0.18\textwidth}{Heuristische Analogie} \\[0.5em]
\parbox[t]{0.14\textwidth}{Kooperativ} & \parbox[t]{0.25\textwidth}{Positive lokale Antwort $c_1>0$} & \parbox[t]{0.32\textwidth}{Spinartige Wechselwirkung, Vielteilchen-Rückkopplung, Korrelationsfunktion} & \parbox[t]{0.18\textwidth}{Lokale formale Korrespondenz} \\[0.5em]
\parbox[t]{0.14\textwidth}{Renormierung} & \parbox[t]{0.25\textwidth}{Assoziativität aufeinanderfolgender Skalenabbildungen, falls Gl.~\eqref{eq:rg_response_homomorphism} gilt} & \parbox[t]{0.32\textwidth}{Projektion vom Theorieraum auf die skalare Antwort eines externen QG-Programms} & \parbox[t]{0.18\textwidth}{Konditionales strukturelles Gate} \\[0.5em]
\parbox[t]{0.14\textwidth}{Temperatur} & \parbox[t]{0.25\textwidth}{Ordnung einer gewählten eindimensionalen Kontrolle nach Umparametrisierung} & \parbox[t]{0.32\textwidth}{Einheiten, thermisches Ensemble, Konjugiertheit oder physikalische Identifikation mit $a$} & \parbox[t]{0.18\textwidth}{Heuristische Umparametrisierung} \\[0.5em]
\parbox[t]{0.14\textwidth}{Kapazität} & \parbox[t]{0.25\textwidth}{Endlicher Antwortendpunkt unter $M_s\leftrightarrow\Phi_0$} & \parbox[t]{0.32\textwidth}{Entropie-, Informations-, Code- oder holographische Kapazität} & \parbox[t]{0.18\textwidth}{Formale Schranke; physikalische Lesart heuristisch} \\
\end{tabular}
\end{ruledtabular}
\end{table*}

Das Mean-Field-Ising-Gleichgewicht erfüllt eine Selbstkonsistenzgleichung der Form $m=\tanh[\beta(B+Jqm)]$; es definiert kein binäres Abel-Kompositionsgesetz für Magnetisierungen. Das Auftreten von $\tanh$ ist daher für sich kein Homomorphismus. Die exakte zweidimensionale Magnetisierung zeigt zudem nicht-Mean-Field-artiges kritisches Verhalten \cite{Yang1952}, und RG-Universalität betrifft Fixpunkte und Skalierungsdaten statt einer gemeinsamen Sigmoidform \cite{WilsonKogut1974}. Tabelle~\ref{tab:metaphor_invariants} führt deshalb nur die explizite normalisierte Antwort-/Normalformabbildung als strukturell; Ferromagnetismus, Temperatur und mikroskopische Kooperation bleiben heuristisch, solange ihre fehlenden Quell-zu-Ziel-Abbildungen nicht hergeleitet sind.

Nachdem wir das ontologische Rahmenwerk etabliert haben, wenden wir uns nun seinen empirischen Konsequenzen zu.

% ===================================================================
% 7. BEOBACHTUNGSMÄSSIGE KONSEQUENZEN
% ===================================================================
\section{Beobachtungsmäßige Konsequenzen und offene Fragen}
\label{sec:observational}

\subsection{Was das Theorem vorhersagt}

Das Saturation Theorem hat zwei Klassen von Testzielen. Tabelle~\ref{tab:observational_ledger} trennt theoremnahe Tests von CRM-spezifischer oder UV-abhängiger Evidenz.

\begin{table*}[t]
\caption{Beobachtungs- und UV-Evidenzledger. Die Tabelle trennt, was durch das Normalformtheorem geprüft wird, von dem, was zum CRM-Fit oder zu einer künftigen UV-Vervollständigung gehört.}
\label{tab:observational_ledger}
\small
\begin{ruledtabular}
\begin{tabular}{@{}llll@{}}
\parbox[t]{0.14\textwidth}{\textbf{Ziel}} & \parbox[t]{0.24\textwidth}{\textbf{Aktuelle Rolle}} & \parbox[t]{0.25\textwidth}{\textbf{Claim-Status}} & \parbox[t]{0.28\textwidth}{\textbf{Nächstes Gate}} \\
\parbox[t]{0.14\textwidth}{Profilfit} & \parbox[t]{0.24\textwidth}{Verträglichkeit der Spätzeitexpansion mit einer $\tanh$-artigen Vorlage} & \parbox[t]{0.25\textwidth}{Nur diagnostisch; CRM wurde mit $\tanh$-Profil konstruiert und ist kein unabhängiger Theoremtest} & \parbox[t]{0.28\textwidth}{Modellunabhängiger Template-Fit ohne CRM-spezifische Priors} \\[0.6em]
\parbox[t]{0.14\textwidth}{Komposition} & \parbox[t]{0.24\textwidth}{Direkter Inhalt der projektiven Normalform} & \parbox[t]{0.25\textwidth}{Theoremnahes empirisches Ziel nach Annahme von D'} & \parbox[t]{0.28\textwidth}{Sättigungsantworten über Rotverschiebungsbins gegen $\mathcal{C}(x,y)=(x+y)/(1+xy)$ prüfen} \\[0.6em]
\parbox[t]{0.14\textwidth}{D'-Defekt} & \parbox[t]{0.24\textwidth}{UV-Auswahl des projektiven Kragens} & \parbox[t]{0.25\textwidth}{Offene physikalische Herleitung, nicht durch die QG-Programmtabelle etabliert} & \parbox[t]{0.28\textwidth}{$\Delta_{D'}=\log R(C)-\log R(x)-\log R(y)$ in einem UV-Modell berechnen oder begrenzen} \\[0.6em]
\parbox[t]{0.14\textwidth}{Parameterwerte} & \parbox[t]{0.24\textwidth}{$k$, $\Phi_0$, $\gamma$ und Übergangsschärfe} & \parbox[t]{0.25\textwidth}{CRM-/UV-spezifisch, nicht vom Theorem bestimmt} & \parbox[t]{0.28\textwidth}{Parameter aus einem mikroskopischen Modell herleiten und gegen CRM-Fits vergleichen} \\[0.6em]
\parbox[t]{0.14\textwidth}{Oszillatorische Residuen} & \parbox[t]{0.24\textwidth}{Mögliche UV-sensitive Spätzeit-Residualstruktur} & \parbox[t]{0.25\textwidth}{Nur Kandidatensignatur; für sich kein Beweis des Quantenursprungs} & \parbox[t]{0.28\textwidth}{Frequenzfenster, Dämpfung, Nullkontrollen und Suchfenster vorregistrieren} \\
\end{tabular}
\end{ruledtabular}
\end{table*}

\textbf{Projektive Randkragen-Vorhersagen} (nach Annahme von D'):
\begin{enumerate}
\item Wenn der projektive Randquotient physikalisch realisiert ist, sollte die kosmologische Expansionsgeschichte mit $\tanh$-ähnlicher Sättigung verträglich sein. Wir betonen, dass das CRM-Ergebnis $\Delta\chi^2 = -5.5$ \cite{Geiger2026b} \textit{keinen} unabhängigen Test des Theorems darstellt, da das CRM mit einem $\tanh$-Profil konstruiert wurde, bevor die Axiome formuliert wurden. Der empirische Gehalt des Theorems ist das \textit{Kompositionsgesetz} $\mathcal{C}(x,y) = (x+y)/(1+xy)$, das eine Vorhersage jenseits jeder individuellen Profilanpassung ist. Ein unabhängiger Test würde entweder (a)~die Anpassung einer modellunabhängigen $\tanh$-Vorlage an SN+CMB+BAO-Daten ohne CRM-spezifische Priors erfordern, oder (b)~das Kompositionsgesetz direkt durch den Vergleich von Sättigungsantworten bei verschiedenen Rotverschiebungs-Bins verifizieren und die Konsistenz mit der Geschwindigkeitsadditionsregel prüfen.
\item Innerhalb des CRM-Fits wird die laufende Krümmungskopplung $\beta_{\mathrm{eff}}(a)$ als monotone Übergangsfunktion berichtet ($\beta_{\mathrm{early}} \approx 2.8 \to \beta_{\mathrm{late}} \approx 2.0$ \cite{Geiger2026b}); theoremnah ist dies ein Konsistenzcheck für Axiom~C, keine unabhängige Herleitung.
\item Kein alternatives Relaxationsprofil (logistisch, Potenzgesetz, etc.) kann Unterstützung aus diesem Theorem beanspruchen, sofern es nicht ebenfalls ein skalenkohärentes Kompositionsgesetz realisiert und den projektiven Quotienten erklärt.
\end{enumerate}

\textbf{UV-abhängige Vorhersagen} (zwischen Vervollständigungen diskriminierend):
\begin{enumerate}
\item Die spezifischen Werte von $k$ (Sättigungsrate), $\Phi_0$ (Sättigungsamplitude) und $\gamma$ (Kopplungskonstante) hängen von der UV-Theorie ab.
\item Die Übergangsschärfe ($n$ im sigmoidalen Übergang) würde eine Universalitätsklasse codieren, sofern eine UV-zu-Kosmologie-Abbildung vorliegt.
\item Sub-Planck-Korrektionen zum $\tanh$-Profil sind als UV-theoriespezifische Kandidaten zu behandeln, solange keine passenden Kontrollen definiert sind.
\end{enumerate}

\subsection{Das fehlende Stück: warum \textit{diese} Parameter?}

Das Saturation Theorem bestimmt die Abel-/Randkragenstruktur und unter D' die projektive \textit{Form}, aber nicht die \textit{Koeffizienten}. Die verbleibende offene Frage ist:

\begin{quote}
\textit{Warum hat das Universum die spezifischen Werte $k \approx 9.8$, $\Phi_0 \approx 0.43$, $\gamma \sim \mathcal{O}(1)\,H_0^{-2}$?}
\end{quote}

Diese Frage ist analog zum Problem der Feinstrukturkonstante: die \textit{Existenz} der elektromagnetischen Kopplung wird durch Eichsymmetrie erzwungen, aber der \textit{Wert} $\alpha \approx 1/137$ wird nicht allein durch die Symmetrie bestimmt.

Mehrere UV-Vervollständigungen machen spezifische Vorhersagen:
\begin{itemize}
\item \textbf{LQC:} $\Phi_0$ wird auf den Barbero-Immirzi-Parameter $\gamma_{\mathrm{BI}}$ abgebildet, der durch Schwarzloch-Entropie eingeschränkt wird \cite{Meissner2004}. Dies würde einen spezifischen Wert von $\Phi_0$ vorhersagen.
\item \textbf{Asymptotische Sicherheit:} Die kritischen Exponenten am UV-Fixpunkt bestimmen die Übergangschärfe $n$, was eine spezifische Vorhersage für die Form von $\beta_{\mathrm{eff}}(a)$ liefert \cite{Reuter1998}.
\item \textbf{Holographie:} Die de Sitter-Entropie $S_{\mathrm{dS}} = \pi/(\Lambda G)$ verbindet $\Phi_0$ mit dem Verhältnis von UV- und IR-Skalen.
\end{itemize}

Ein konkreter Testfall ist die $\sqrt{\pi}$-Vermutung von Paper~II \cite{Geiger2026b}: das CRM sagt einen MOND-ähnlichen Verstärkungsfaktor $\mu_{\mathrm{eff}} = \sqrt{\pi}$ voraus, der aus der Gauß-Normalisierung der Sättigungs-Partitionsfunktion hervorgeht. Die Ableitung dieses Wertes aus einer spezifischen UV-Vervollständigung würde direkt die Quantengravitations-Parameter mit der beobachteten Galaxien-Dynamik verbinden.

\subsection{Skalarfeld-Oszillationen als QG-Signatur}

Paper~III \cite{Geiger2026c} zeigt, dass das Pöschl-Teller-Potential ein diskretes Spektrum von Bindungszuständen unterstützt. Im kosmologischen Kontext entsprechen diese oszillatorischen Korrektionen zur Expansionsrate in späten Zeiten:
\begin{equation}
H^2(a) = H_{\mathrm{smooth}}^2(a)\left[1 + \epsilon \cdot e^{-\Gamma a} \cos(\omega a + \delta)\right]\,,
\end{equation}
mit $\epsilon \ll 1$. Diese Oszillationen wären eine Kandidatensignatur mikroskopischer Struktur im Sättigungsmechanismus -- der Analogie nach verwandt mit diskreten Energieniveaus eines Quantensystems. Eine Detektion in hochpräzisen BAO- oder SN-Daten würde nur dann einen UV-sensitiven Residualkanal stützen, wenn Frequenzfenster, Dämpfungsskala und Nullkontrollen vor der Anpassung festgelegt wurden; sie wäre für sich kein Beweis des Quantenursprungs der CRM-Sättigung.


% ===================================================================
% 8. DISKUSSION
% ===================================================================
\section{Diskussion und Ausblick}
\label{sec:discussion}

\subsection{Was wurde erreicht}

\begin{enumerate}
\item \textbf{Eine konditionale Normalform-/No-Go-Aussage:} Jede kosmologische Theorie, die die Axiome~A--D, B', $D^\pm$ und die projektive Doppelrand-Quotienten-Annahme D' erfüllt, muss eine $\tanh$-artige Sättigung produzieren. Ohne D' liefern dieselben strukturellen Voraussetzungen die Abel-/Randkragenstruktur.

\item \textbf{Universalität eingegrenzt:} Das Theorem definiert eine konditionale Abel-/Randkragen-Äquivalenzklasse für Antworten, die seine Gates erfüllen. Die QG-Beispiele liefern derzeit Motive und Abbildungspflichten, keine Wilsonsche Universalitätsklasse oder einen gemeinsamen mikroskopischen Mechanismus; der projektive $\tanh$-Repräsentant hängt zusätzlich von D' ab.

\item \textbf{CRM kalibriert:} Das CRM's Sättigungs-ODE $dX/da = k(1-X^2)$ kann unter den genannten Axiomen als projektive Normalform kapazitätsbegrenzter kooperativer Relaxation gelesen werden, nicht nur als phänomenologisches Fitting.

\item \textbf{Ontologische Interpretation:} In diesem Rahmen kann Quantengravitation nicht nur als Metrikquantisierung, sondern auch als Theorie geometrischer Ordnung gelesen werden -- als Übergang von einem prä-geometrischen Quantensystem zu klassischer Raumzeit.
\end{enumerate}

\subsection{Relation zum CRM-Programm}

Dieses Paper liefert die bedingte axiomatische Schicht des CRM-Programms, vervollständigt aber nicht dessen physikalischen Bogen. Nach dem Wahrheitsaudit und der Route-3-Rekonstruktion vom Juli~2026 gilt folgender logischer Status:
\begin{enumerate}
\item Die \textit{Form} der Dynamik ($\tanh$) ist ein mathematisches Theorem, sobald die signierten inneren Kompositionsvoraussetzungen und der projektive Randquotient D' akzeptiert werden; ohne D' beweist das Theorem die Abel-/Randkragenstruktur.
\item Die historische direkte Spurkopplungsroute aus Paper~III ist ein Fehlaudit; Route~3 ist ein Wirkungskandidat, dessen vollständige FLRW-, Hamilton-, PPN-, Lensing- und Datengates offen bleiben.
\item Die berichteten Fitparameter beschränken einen effektiven phänomenologischen Proxy, noch keine validierte Route-3- oder $f(R)$-Realisierung.
\item Die \textit{UV-Vervollständigung} bleibt offen -- und muss erklären, warum der projektive Randquotient und nicht ein anderer glatter Randkragen die physikalische Kompositionsvariable ist.
\end{enumerate}

\subsection{Post-Zookeeper-Roadmap für Stabilitätserweiterungen}

Das Post-Zookeeper-Programm ist der vom Autor vorgeschlagene künftige Rahmen zur Klassifikation zulässiger UV-Antwortkanäle; es wird im Beweis des Theorems nicht verwendet.

\subsection{Limitationen und Einschränkungen}

\begin{enumerate}
\item \textbf{Axiom D und D' sind stark.} Das Renormalisierungsgruppen-Kompositionsgesetz ist restriktiv, und die projektive Doppelrand-Quotienten-Annahme ist eine zusätzliche mathematische und physikalische Wahl. Theorien ohne gut definiertes RG-Fließverhalten (einige diskrete Ansätze, Multiversum-Szenarien) erfüllen möglicherweise D nicht; Theorien mit einem anderen Randkragen können D erfüllen, aber D' verfehlen. Diese Unterscheidung ist wesentlich: D liefert Skalenkohärenz, D' wählt den $\tanh$-Repräsentanten.

\item \textbf{Der Beweis ist algebraisch, nicht konstruktiv.} Das Theorem stellt die Abel-Struktur der Antwort und unter D' ihren projektiven $\tanh$-Repräsentanten fest, konstruiert aber die UV-Theorie nicht. Die ``Reparametrisierung'' $u(g)$ absorbiert mikroskopische Informationen; $u(g)$ aus ersten Prinzipien zu bestimmen erfordert das Lösen der spezifischen QG-Theorie. Ein potentieller Einwand ist, dass jede beschränkte monotone Funktion trivial als $\tanh(\alpha\,u(g))$ für ein geeignetes $u(g)$ geschrieben werden kann, was die Profilform inhaltsleer macht. Der nichttriviale Inhalt liegt daher im \textit{Kompositionsgesetz} und für die $\tanh$-Auswahl im projektiven Quotientengesetz $R(C)=R(x)R(y)$.

\item \textbf{Kosmologische Einschränkung und Theoriedrift.} Die Axiome sind für einen ausgewählten kosmologischen Antwortkanal formuliert und begründen keine Physik Schwarzer Löcher, Gravitationswellen, Screening- oder Clusterphänomenologie. Frühere Chamäleon- und Bullet-Cluster-Claims gehören zur abgelösten Route und folgen nicht aus diesem Theorem. In der aktiven Route-3-Rekonstruktion ist $c_T=1$ nur auf der transversal-spurlosen Familie bestätigt; Skalar-Hamiltonstabilität, FLRW-Antwort, Lensing und Screening bleiben getrennte Gates.

\item \textbf{Die kontinuierliche Gruppenklassifikation und Einzelfeld-Einschränkung.} Der Beweis beruht auf Aczél's Darstellungstheorem für kontinuierliche assoziative Operationen \cite{Aczel1966} und der Klassifikation eindimensionaler formaler Gruppengesetze \cite{Hazewinkel1978}. Während diese Resultate gut etabliert sind, erfordert die physikalische Anwendbarkeit die Annahme, dass die Komposition \textit{eindimensional} ist -- d.h. dass es eine einzige relevante Richtung gibt. Dies ist eine echte Limitation: das CRM selbst (Paper~IV \cite{Geiger2026d}) führt einen Vektorsektor neben dem Skalar ein, und realistische modifizierte Gravitationstheorien (z.B. Horndeski, beyond-Horndeski) haben generisch mehrere gekoppelte Freiheitsgrade. Jedoch merken wir an, dass das Theorem auf den \textit{Skalarsektor in Isolation} angewendet wird: der CRM's Skalar-Sättigungs-Kanal $\Omega_\Phi(a) = \Phi_0 \tanh(k(a - a_{\mathrm{trans}}))$ ist eine eindimensionale Antwort, die die Axiome unabhängig vom Vektorsektor erfüllt. Die Mehrfeld-Erweiterung (Klassifikation mehrdimensionaler kommutativer Gruppengesetze, einschließlich elliptischer und Lubin-Tate formaler Gruppen) ist ein offenes mathematisches Problem, das das Theorem verstärken würde, wird aber für das Skalarsektor-Resultat nicht benötigt.

\item \textbf{Randkragen-Freiheit.} Glattes Randverhalten allein wählt $\mathrm{arctanh}$ nicht aus. Der Gegen-Kragen $\phi_\varepsilon=\operatorname{arctanh}+\varepsilon x$ erhält einen glatten positiven Randkragen, ändert aber die Abel-Koordinate. Die Eindeutigkeitsauswahl wird daher bewusst auf D' verlagert: Sobald der projektive Quotient $R=(1+x)/(1-x)$ unter Komposition multiplizieren muss, folgt $\mathrm{arctanh}=\frac12\log R$ unmittelbar und bis auf Skalierung eindeutig.
\end{enumerate}

\subsection{Das No-Go-Theorem als klasseninterner Filter}

Innerhalb des Rahmens A--D+B'+$D^\pm$ bedeutet ein Versagen von D', dass das $\tanh$-CRM-Gesetz nicht als projektive Normalform ausgewählt wird. Ein solches Versagen falsifiziert das UV-Modell nicht; es verortet es außerhalb der projektiven CRM-Unterklasse oder identifiziert einen messbaren D'-Defekt. Theorien, die die Axiome A--D zusammen mit B' und $D^\pm$ erfüllen, sind nur auf ein Abel-/Randkragen-Sättigungsgesetz eingeschränkt. Theorien, die diese strukturellen Voraussetzungen und D' erfüllen, sind auf $\tanh$-Sättigung in ihrem kosmologischen Antwortkanal eingeschränkt.

\subsection{Relation zu anderen Rahmenwerken}

\textit{Swampland-Vermutungen.} Die de~Sitter-Swampland-Vermutung \cite{Obied2018} postuliert $|\nabla V|/V \geq c \sim \mathcal{O}(1)$ in jeder Quantengravitations-Landschaft und benachteiligt eine strikte kosmologische Konstante. Das Saturation Theorem ist komplementär: es adressiert nicht die Existenz von de~Sitter-Vakua, aber schränkt den \textit{dynamischen Zugang} zum Spätzeitattraktor ein. Eine CRM-artige Sättigung mit $\Phi_0 < \infty$ ist kompatibel mit der Swampland-Schranke, da die effektive Dunkle-Energie-Dichte $\rho_\Phi(a)$ dynamisch anstatt als festes $\Lambda$ entwickelt.

\textit{EFT der Dunklen Energie.} Die effektive Feldtheorie der Dunklen Energie \cite{Gubitosi2013} parametrisiert Abweichungen vom $\Lambda$CDM auf der Ebene der gestörten Wirkung. Das Saturation Theorem operiert auf einer anderen Ebene: es schränkt die \textit{Hintergrund}-Antwortfunktion ein, nicht die perturbative Operatorbasis. Eine zukünftige Aufgabe ist, die $\tanh$-Sättigung in der EFT-of-DE-Sprache auszudrücken und die CRM-Parameter $(k, \Phi_0, \gamma)$ auf die EFT-Operatorkoeffizienten $\{\alpha_i(t)\}$ abzubilden.

\subsection{Ausblick}

Drei Richtungen sind unmittelbar:

\begin{enumerate}
\item \textbf{Parameterherleitung.} Der nächste Schritt ist, $k$, $\Phi_0$ und $\gamma$ aus einer spezifischen UV-Vervollständigung herzuleiten. Der vielversprechendste Kandidat ist die LQG/holographische Schnittmenge, wo der Barbero-Immirzi-Parameter und die holographische Entropie-Schranke gemeinsam die Sättigungsparameter fixieren könnten.

\item \textbf{Mehrfeld-Erweiterung.} Die Erweiterung des Theorems auf Theorien mit mehreren Sättigungs-Kanälen (z.B. Skalar + Vektor im CRM \cite{Geiger2026d}) erfordert die Klassifikation mehrdimensionaler kommutativer Gruppengesetze (das höherdimensionale Analogon der Abel-Gleichungs-Reduktion). Dies ist mathematisch reichhaltiger und könnte Beschränkungen auf die Vektorsektor-Parameter $K_B$ und $\mathcal{F}$ liefern.

\item \textbf{Schwarzloch-Inneres.} Wenn das Saturation Theorem auf Schwarzloch-Innere (mit der Radialkoordinate die Skalenfaktor ersetzend) angewendet wird und der projektive Randquotient dort der richtige innere Kragen ist, sagt es eine $\tanh$-artige Singularitätsauflösung voraus -- konsistent mit LQG-Bounce-Modellen \cite{Ashtekar2006} und dem Planck-Star-Szenario \cite{Rovelli2014}.
\end{enumerate}


% ===================================================================
% 9. SCHLUSSFOLGERUNG
% ===================================================================
\section{Schlussfolgerung}
\label{sec:conclusion}

Wir haben ein korrigiertes Saturation Theorem bewiesen: endliche geometrische Kapazität, kooperative Verstärkung, monotone kosmische Kontrolle und Renormalisierungsgruppen-Komposition reduzieren eindimensionale kosmologische Sättigung auf eine Abel-Koordinaten-/Randkragenstruktur. Das $\tanh$-Profil folgt, wenn der Sättigungsrand zusätzlich durch den projektiven Quotienten $R=(1+x)/(1-x)$ normalisiert wird.

Das Theorem hat drei Konsequenzen:

\begin{enumerate}
\item Das CRM's Sättigungs-ODE $dX/da = k(1-X^2)$ kann als \textit{projektive Normalform} von kapazitätsbegrenzter kooperativer Relaxation gelesen werden, sobald der Randquotient fixiert ist.

\item Motive aus LQG, Asymptotic Safety, Stringtheorie, Causal Sets und Noncommutative Geometry können mit den Abel-/Randkragen-Gates verglichen werden, doch ein externer Quell-zu-Antwort-Homomorphismus wird hier nicht hergeleitet. Der konkrete $\tanh$-Repräsentant erfordert zusätzlich den projektiven Quotienten.

\item Das Paper bietet einen Interpretationsrahmen, in dem Quantengravitation nicht nur als Metrikquantisierung, sondern auch als Theorie geometrischer Ordnung untersucht wird -- als Übergang von einem prä-geometrischen Quantensystem zu klassischer Raumzeit.
\end{enumerate}

Das verbleibende offene Problem -- die spezifischen Parameter $(k, \Phi_0, \gamma)$ und den projektiven Randquotienten aus einer UV-Vervollständigung herzuleiten -- ist das Analogon zur Ableitung der Feinstrukturkonstante aus einer Theorie der Quantenelektrodynamik. Das Theorem schränkt die \textit{Normalform} jeder Theorie ein, die die genannten Hypothesen erfüllt, sobald der Randquotient fixiert ist.

\begin{quote}
\textit{Das Sättigungsprofil ist durch Skalenkohärenz plus projektiven Randquotienten auf den $\tanh$-Typ eingeschränkt; bloße glatte Randregularität liefert die breitere Abel-Kragenklasse.}
\end{quote}


% ===================================================================
% KI-OFFENLEGUNG
% ===================================================================
\begin{acknowledgments}
Diese Arbeit wurde als unabhängige Forschung ohne institutionelle Finanzierung durchgeführt. Öffentliche wissenschaftliche Recheninfrastruktur wurde im gesamten CRM-Programm genutzt.

\textit{Werkzeugeinsatz.} KI-basierte Werkzeuge unterstützten den Arbeitsprozess nicht-autorschaftlich bei Literaturorganisation, Entwurfs- und Codeassistenz, sprachlicher Überarbeitung, Konsistenzprüfungen sowie der \LaTeX-/PDF-Qualitätssicherung. Quellenauswahl, mathematische Argumente, physikalische Hypothesen, axiomatische Entscheidungen, wissenschaftliche Interpretation und finale Entscheidungen wurden vom Autor geprüft und verantwortet. Die Werkzeuge sind weder Autoren noch Koautoren, werden nicht als Beitragende gewürdigt und dienten weder als Wahrheitsquelle noch als unabhängige Validierungsinstanz. Der Autor trägt die volle Verantwortung für den Inhalt.

\textit{Finanzierungsinformationen.} Diese Forschung erhielt keine externe Finanzierung.

\textit{Datenverfügbarkeit.} Dieses Paper ist rein theoretisch. Das öffentliche CRM-Repository mit Analyseskripten, Datenprodukten und Paper-Quelldateien ist verfügbar unter \url{https://github.com/research-line/crm-cosmology}; der Zenodo-Concept-DOI für CRM I--IV ist \href{https://doi.org/10.5281/zenodo.18728935}{10.5281/zenodo.18728935}.
\end{acknowledgments}


% ===================================================================

% ===================================================================
% BIBLIOGRAPHY
% ===================================================================
\begin{thebibliography}{50}

\bibitem{Geiger2026}
L.~Geiger (2026).
Game-Theoretic Cosmology and the Curvature Relaxation Model: Nash Equilibria Between Null Space and Spacetime Bubble.
Companion paper (Paper~I).
\url{https://github.com/research-line/crm-cosmology/tree/main/papers}.
\href{https://doi.org/10.5281/zenodo.18728935}{doi:10.5281/zenodo.18728935}.

\bibitem{Geiger2026b}
L.~Geiger (2026).
Eliminating the Dark Sector: Unifying the Curvature Relaxation Model with MOND.
Companion paper (Paper~II).
\href{https://doi.org/10.5281/zenodo.18728935}{doi:10.5281/zenodo.18728935}.

\bibitem{Geiger2026c}
L.~Geiger (2026).
Microscopic Foundations of the Curvature Relaxation Model: From Quantum Geometry to Macroscopic Saturation.
Companion paper (Paper~III).
\href{https://doi.org/10.5281/zenodo.18728935}{doi:10.5281/zenodo.18728935}.

\bibitem{Geiger2026d}
L.~Geiger (2026).
The Galactic--Cosmological Nexus: Connecting CRM to MOND Dynamics via Curvature Saturation.
Companion paper (Paper~IV).
\href{https://doi.org/10.5281/zenodo.18728935}{doi:10.5281/zenodo.18728935}.

\bibitem{Donoghue1994}
J.~F.~Donoghue (1994).
General relativity as an effective field theory: The leading quantum corrections.
\textit{Physical Review D}, 50(6), 3874.
DOI: 10.1103/PhysRevD.50.3874.
arXiv:gr-qc/9405057.

\bibitem{Eichhorn2020}
A.~Eichhorn,
\textit{Asymptotically safe gravity},
arXiv:2003.00044 (2020).

\bibitem{EichhornHeld2018}
A.~Eichhorn and A.~Held,
\textit{Top mass from asymptotic safety},
Phys.\ Lett.\ B \textbf{777}, 217--221 (2018).
DOI: 10.1016/j.physletb.2017.12.040.
arXiv:1707.01107.

\bibitem{Weinberg1979}
S.~Weinberg (1979).
Ultraviolet divergences in quantum theories of gravitation.
In S.~W.~Hawking \& W.~Israel (Eds.), \textit{General Relativity: An Einstein Centenary Survey}, pp.\ 790--831.
Cambridge University Press.

\bibitem{Reuter1998}
M.~Reuter (1998).
Nonperturbative evolution equation for quantum gravity.
\textit{Physical Review D}, 57(2), 971.
DOI: 10.1103/PhysRevD.57.971.
arXiv:hep-th/9605030.

\bibitem{Percacci2017}
R.~Percacci (2017).
\textit{An Introduction to Covariant Quantum Gravity and Asymptotic Safety}.
World Scientific.

\bibitem{Polchinski1998}
J.~Polchinski (1998).
\textit{String Theory}, Vols.~1--2.
Cambridge University Press.

\bibitem{Rovelli2004}
C.~Rovelli (2004).
\textit{Quantum Gravity}.
Cambridge University Press.

\bibitem{Thiemann2007}
T.~Thiemann (2007).
\textit{Modern Canonical Quantum General Relativity}.
Cambridge University Press.

\bibitem{Maldacena1999}
J.~Maldacena (1998).
The large-$N$ limit of superconformal field theories and supergravity.
\textit{Adv.\ Theor.\ Math.\ Phys.}\ \textbf{2}, 231--252.
[Reprinted in \textit{Int.\ J.\ Theor.\ Phys.}\ \textbf{38}, 1113--1133 (1999).]
DOI: 10.4310/ATMP.1998.v2.n2.a1.
arXiv:hep-th/9711200.

\bibitem{Sorkin2003}
R.~D.~Sorkin (2005).
Causal sets: Discrete gravity.
In \textit{Lectures on Quantum Gravity}, pp.\ 305--327.
Springer.
DOI: 10.1007/0-387-24992-3\_7.
arXiv:gr-qc/0309009.

\bibitem{Surya2019}
S.~Surya (2019).
The causal set approach to quantum gravity.
\textit{Living Reviews in Relativity}, 22, 5.
DOI: 10.1007/s41114-019-0023-1.

\bibitem{Connes1994}
A.~Connes (1994).
\textit{Noncommutative Geometry}.
Academic Press.

\bibitem{Bekenstein1981}
J.~D.~Bekenstein (1981).
Universal upper bound on the entropy-to-energy ratio for bounded systems.
\textit{Physical Review D}, 23(2), 287--298.
DOI: 10.1103/PhysRevD.23.287.

\bibitem{Almheiri2015}
A.~Almheiri, X.~Dong, \& D.~Harlow (2015).
Bulk locality and quantum error correction in AdS/CFT.
\textit{Journal of High Energy Physics}, 2015, 163.
DOI: 10.1007/JHEP04(2015)163.

\bibitem{Aczel1966}
J.~Acz\'el (1966).
\textit{Lectures on Functional Equations and Their Applications}.
Academic Press, New York.

\bibitem{Hazewinkel1978}
M.~Hazewinkel (1978).
\textit{Formal Groups and Applications}.
Academic Press.

\bibitem{Ashtekar2006}
A.~Ashtekar, T.~Pawlowski, \& P.~Singh (2006).
Quantum nature of the Big Bang.
\textit{Physical Review Letters}, 96(14), 141301.
DOI: 10.1103/PhysRevLett.96.141301.
arXiv:gr-qc/0602086.

\bibitem{Bojowald2008}
M.~Bojowald (2008).
Loop quantum cosmology.
\textit{Living Reviews in Relativity}, 11, 4.
DOI: 10.12942/lrr-2008-4.

\bibitem{Bonanno2002}
A.~Bonanno \& M.~Reuter (2002).
Cosmology of the Planck era from a renormalization group for quantum gravity.
\textit{Physical Review D}, 65(4), 043508.
DOI: 10.1103/PhysRevD.65.043508.
arXiv:hep-th/0106133.

\bibitem{Ryu2006}
S.~Ryu \& T.~Takayanagi (2006).
Holographic derivation of entanglement entropy from the anti-de~Sitter space/conformal field theory correspondence.
\textit{Physical Review Letters}, 96(18), 181602.
DOI: 10.1103/PhysRevLett.96.181602.
arXiv:hep-th/0603001.

\bibitem{deBoer2000}
J.~de Boer, E.~P.~Verlinde, \& H.~L.~Verlinde (2000).
On the holographic renormalization group.
\textit{Journal of High Energy Physics}, 2000(08), 003.
DOI: 10.1088/1126-6708/2000/08/003.
arXiv:hep-th/9912012.


\bibitem{Yang1952}
C.~N.~Yang (1952).
The spontaneous magnetization of a two-dimensional Ising model.
\textit{Physical Review}, 85, 808--816.
DOI: 10.1103/PhysRev.85.808.

\bibitem{WilsonKogut1974}
K.~G.~Wilson \& J.~Kogut (1974).
The renormalization group and the $\epsilon$ expansion.
\textit{Physics Reports}, 12(2), 75--199.
DOI: 10.1016/0370-1573(74)90023-4.

\bibitem{Meissner2004}
K.~A.~Meissner (2004).
Black-hole entropy in Loop Quantum Gravity.
\textit{Classical and Quantum Gravity}, 21(22), 5245.
DOI: 10.1088/0264-9381/21/22/015.
arXiv:gr-qc/0407052.

\bibitem{Rovelli2014}
C.~Rovelli \& F.~Vidotto (2014).
Planck stars.
\textit{International Journal of Modern Physics D}, 23(12), 1442026.
DOI: 10.1142/S0218271814420267.
arXiv:1401.6562.

\bibitem{PDG2024}
S.~Navas et~al. (Particle Data Group) (2024).
Review of Particle Physics.
\textit{Physical Review D}, 110, 030001.
DOI: 10.1103/PhysRevD.110.030001.

\bibitem{Obied2018}
G.~Obied, H.~Ooguri, L.~Spodyneiko, \& C.~Vafa (2018).
De Sitter Space and the Swampland.
arXiv:1806.08362.

\bibitem{Gubitosi2013}
G.~Gubitosi, F.~Piazza, \& F.~Vernizzi (2013).
The Effective Field Theory of Dark Energy.
\textit{Journal of Cosmology and Astroparticle Physics}, 2013(02), 032.
DOI: 10.1088/1475-7516/2013/02/032.
arXiv:1210.0201.

\bibitem{Ambjorn2005}
J.~Ambj\o{}rn, J.~Jurkiewicz, \& R.~Loll (2005).
The spectral dimension of the universe is scale dependent.
\textit{Physical Review Letters}, 95(17), 171301.
DOI: 10.1103/PhysRevLett.95.171301.
arXiv:hep-th/0505113.

\bibitem{Loll2020}
R.~Loll (2020).
Quantum gravity from causal dynamical triangulations: a review.
\textit{Classical and Quantum Gravity}, 37(1), 013002.
DOI: 10.1088/1361-6382/ab57c7.
arXiv:1905.08669.

\bibitem{HartnollYang2025}
S.~A.~Hartnoll \& M.~Yang (2025).
The conformal primon gas at the end of time.
\textit{Journal of High Energy Physics}, 2025(07), 281.
DOI: 10.1007/JHEP07(2025)281.
arXiv:2502.02661.

\bibitem{Julia1990}
B.~L.~Julia (1990).
Statistical theory of numbers.
In J.~M.~Luck, P.~Moussa, \& M.~Waldschmidt (Eds.), \textit{Number Theory and Physics},
Springer Proceedings in Physics, Vol.~47, pp.\ 276--293.
Springer-Verlag, Berlin.

\bibitem{PapadopoulosTroyanov2008}
A.~Papadopoulos \& M.~Troyanov (2008).
Harmonic symmetrization of convex sets and of Finsler structures, with applications to Hilbert geometry.
arXiv:0807.0335.

\bibitem{RainioVuorinen2023}
O.~Rainio \& M.~Vuorinen (2023).
Hilbert metric in the unit ball.
arXiv:2303.03753.

\bibitem{Selberg1956}
A.~Selberg (1956).
Harmonic analysis and discontinuous groups in weakly symmetric Riemannian spaces with applications to Dirichlet series.
\textit{Journal of the Indian Mathematical Society}, 20, 47--87.
DOI: 10.18311/JIMS/1956/16985.

\bibitem{Ungar2007}
A.~A.~Ungar (2007).
Einstein's velocity addition law and its hyperbolic geometry.
\textit{Computers \& Mathematics with Applications}, 53(8), 1228--1250.
DOI: 10.1016/j.camwa.2006.05.028.

\end{thebibliography}

\end{document}
